% Permission to make digital/hard copy of part or all of this work for
% personal or classroom use is granted without fee provided that copies
% are not made or distributed for profit or commercial advantage, the
% copyright notice, the title of the publication and its date appear, and
% notice is given that copying is by permission of the authors. To copy
% otherwise, to republish, to post on servers, or to redistribute to lists
% requires prior specific permission and/or a fee.
%
% Copyright (C) 1998 by Gernot Heiser, The University of New South Wales.


\documentclass[a4paper,11pt,twoside,dvips]{report}
\usepackage{Bkit}
\usepackage{subfigure}
\usepackage{epsfig}
\usepackage{rcs}
\usepackage{times}
\usepackage{cite}
\usepackage{changebar}

\RCS$Revision: 1.16 $
\RCS$Date: 1999/08/11 02:33:34 $

\sloppy

\title{\Huge L4 User Manual\\[20pt]
       \LARGE Version \RCSRevision}

\newcommand{\Tilde}{\raisebox{-0.6ex}{\~{}}}
\author{Alan Au, Gernot Heiser\\
       {\normalsize School of Computer Science and Engineering}\\
       {\normalsize The University of New South Wales}\\
       {\normalsize Sydney 2052, Australia}\\
       {\normalsize \{alanau,gernot\}@cse.unsw.edu.au}}

\date{\RCSDate\\[1ex] Version 1.5 is identical to UNSW-CSE-TR-9801 of June 1998\\
	\vfill\sf
	\epsfig{file=eps/unicrest-screen.eps,height=25mm}\\[2ex]
	School of Computer Science and Engineering\\
	The University of New South Wales\\
	Sydney 2052, Australia \vspace{-10mm}}

\newcounter{example}[chapter]
\newcommand{\example}[1]{\addtocounter{example}{1}
			\underline{\emph{Example \thechapter.\theexample}}
			 - #1 \\}
 
\newcommand{\micro}{$\mu$}
\renewcommand{\textfraction}{0.05}
\renewcommand{\topfraction}{0.95}
\renewcommand{\bottomfraction}{0.95}
\newcommand{\sigz}{\mbox{$\sigma_{0}$}}



\newcommand{\smaller}[1]{{\mathchoice
           {\raisebox{1pt}{$\scriptstyle #1$}}
           {\raisebox{1pt}{$\scriptstyle #1$}}
           {\raisebox{0.5pt}{$\scriptscriptstyle #1$}}
           {\raisebox{0.5pt}{$\scriptscriptstyle #1$}}
}}
\newcommand{\smallerlow}[1]{{\mathchoice
           {\mbox{$\scriptstyle \,#1\,$}}
           {\mbox{$\scriptstyle \,#1\,$}}
           {\mbox{$\scriptscriptstyle \,#1\,$}}
           {\mbox{$\scriptscriptstyle \,#1\,$}}
}}
\newcommand{\IN}{\smaller{\in}}
\newcommand{\LT}{\smaller{<}}
\newcommand{\LE}{\smaller{\le}}
\newcommand{\GT}{\smaller{>}}
\newcommand{\GE}{\smaller{\ge}}
\newcommand{\EQ}{\smaller{=}}
\newcommand{\NE}{\smaller{\neq}}
\newcommand{\CUP}{\smallerlow{\cup}}
\newcommand{\PLUS}{\smaller{+}}

\newcommand{\nil}{--nil--\ }
\newcommand{\never}{--never--\ }
\newcommand{\undef}{$\sim$}

\newcommand{\stdtabs}{\sf\=XX\=XX\=XX\=XX\=XX\=XX\=XX\=Xx\=Xx\=Xx\=Xx\=XXXXX\=\kill\+\+\\}

\newcommand{\alg}[1]
 {{\noindent
 \begin{minipage}{\textwidth}
 \begin{tabbing}\stdtabs\+ #1 \end{tabbing}
 \end{minipage}}}

\newcommand{\asmtabs}{\sf\=XYYXXXX\=XYYXXXX\=XXXXX\=XXXXXXXXXXXX\=\kill\+\+\\[-5pt]}

\newcommand{\asm}[1]
 {{\noindent\tt
 \begin{minipage}{\textwidth}
 \begin{tabbing}\asmtabs #1 \end{tabbing}
 \end{minipage}}}


\newcommand{\WHILE}{{\bf while}\ }
\newcommand{\UNTIL}{{\bf until}\ }
\newcommand{\DO}{{\bf do}\ \+}
\newcommand{\OD}{{\<\bf od}\ \-}
\newcommand{\IF}{{\bf if}\ \+}
\newcommand{\FI}{{\<\<\<\bf fi}\ \-\-\-}
\newcommand{\THEN}{{\bf then}\ \>\>\+\+}
\newcommand{\ELSE}{{\<\<\bf else}\>\>}
\newcommand{\ELIF}{{\<\<\<\bf elif}\ \-\-}

\newcommand{\While}{{\bf while}\ }
\newcommand{\Until}{{\bf until}\ }
\newcommand{\Do}{{\bf do}\ }
\newcommand{\Od}{{\bf od}\ }
\newcommand{\If}{{\bf if}\ }
\newcommand{\Fi}{{\bf fi}\ }
\newcommand{\Then}{{\bf then}\ }


\newcommand{\dd}{D}
\newcommand{\dci}{d_{c_i}}
\newcommand{\Dci}{d_{c_i,\dot{v}}}
\newcommand{\dv}{V}
\newcommand{\ddv}{V'}
\newcommand{\vci}{v_{c_i}}
\newcommand{\cj}{c_j}

\newcommand{\dw}{\mbox{\it write}}
\newcommand{\ndw}{\overline{\mbox{\it write}}}

\newcommand{\allj}{\forall_{j\neq i}:\ \alias\mapsto\cj}

\newcommand{\rdy}{\ \Box}
\newcommand{\sync}{\ \Box}

\newcommand{\high}{\mbox{\rule{0pt}{\baselineskip}}}
\newcommand{\High}{\mbox{$\frac{\high}{\high}$}}

\newlength{\Xlen}
\settowidth{\Xlen}{0}
\newcommand{\X}{\makebox[\Xlen]{}}


\newlength{\Up}\setlength{\Up}{-\baselineskip}
\newlength{\Upp}\setlength{\Upp}{\Up}\addtolength{\Upp}{\Upp}
\newlength{\Upppp}\setlength{\Upppp}{\Upp}\addtolength{\Upppp}{\Upppp}

\newlength{\PPre}
\newlength{\Params}

\newif\ifintel
\intelfalse
\newlength{\scindent}
\setlength{\scindent}{140pt}

\newcommand{\generalsc}[1]{\intelfalse\setlength{\scindent}{#1}}
\newcommand{\intelsc}[1]{\inteltrue\setlength{\scindent}{#1}}

\newcommand{\sccode}[1]{\global\def\syscode{#1}}


\newcommand{\scin}[1]{(#1)\ }
\newcommand{\scout}[1]{{\Large\ $\rightarrow$\ } (#1)\\}





\newcommand{\EAX}[2]{{\em #1}&{\footnotesize EAX}&&{\footnotesize EAX}&{\em #2}\\}
\newcommand{\EBX}[2]{{\em #1}&{\footnotesize EBX}&&{\footnotesize EBX}&{\em #2}\\}
\newcommand{\ECX}[2]{{\em #1}&{\footnotesize ECX}&&{\footnotesize ECX}&{\em #2}\\}
\newcommand{\EDX}[2]{{\em #1}&{\footnotesize EDX}&&{\footnotesize EDX}&{\em #2}\\}
\newcommand{\ESI}[2]{{\em #1}&{\footnotesize ESI}&&{\footnotesize ESI}&{\em #2}\\}
\newcommand{\EDI}[2]{{\em #1}&{\footnotesize EDI}&&{\footnotesize EDI}&{\em #2}\\}
\newcommand{\EBP}[2]{{\em #1}&{\footnotesize EBP}&&{\footnotesize EBP}&{\em #2}\\}
\newcommand{\EAL}[2]{{\em #1}&{\footnotesize  AL}&&{\footnotesize EAX}&{\em #2}\\}
\newcommand{\EAH}[1]{{\em #1}&{\footnotesize  AH}&&&\\}
\newcommand{\EBL}[2]{{\em #1}&{\footnotesize  BL}&&{\footnotesize EBX}&{\em #2}\\}
\newcommand{\BH}[1]{{\em #1}&{\footnotesize  BH}&&&\\}
\newcommand{\ECL}[2]{{\em #1}&{\footnotesize  CL}&&{\footnotesize ECX}&{\em #2}\\}
\newcommand{\CH}[1]{{\em #1}&{\footnotesize  CH}&&&\\}

\newcommand{\EAXX}[3]{{\em #1}&{\footnotesize EAX}&&{\footnotesize EAX}&{\em #2}&~~{\scriptsize /}~~~{\em #3}\\}
\newcommand{\EBXX}[3]{{\em #1}&{\footnotesize EBX}&&{\footnotesize EBX}&{\em #2}&~~{\scriptsize /}~~~{\em #3}\\}
\newcommand{\ECXX}[3]{{\em #1}&{\footnotesize ECX}&&{\footnotesize ECX}&{\em #2}&~~{\scriptsize /}~~~{\em #3}\\}
\newcommand{\EDXX}[3]{{\em #1}&{\footnotesize EDX}&&{\footnotesize EDX}&{\em #2}&~~{\scriptsize /}~~~{\em #3}\\}
\newcommand{\ESIX}[3]{{\em #1}&{\footnotesize ESI}&&{\footnotesize ESI}&{\em #2}&~~{\scriptsize /}~~~{\em #3}\\}
\newcommand{\EDIX}[3]{{\em #1}&{\footnotesize EDI}&&{\footnotesize EDI}&{\em #2}&~~{\scriptsize /}~~~{\em #3}\\}
\newcommand{\EBPX}[3]{{\em #1}&{\footnotesize EBP}&&{\footnotesize EBP}&{\em #2}&~~{\scriptsize /}~~~{\em #3}\\}

\newcommand{\USP}[2]{{\em #2}&%
\multicolumn{5}{l}{{\footnotesize $\langle$SP#1$\rangle$}}\\}

\newcommand{\Or}{\ {\scriptsize /}\ \ }

\newenvironment{param}[1]
{%
\setlength{\Params}{\textwidth}\addtolength{\Params}{-140pt}%
\par\vspace{5pt}%
\noindent\begin{minipage}{\textwidth}
\noindent {\em #1}\\[\Up]%
\noindent\begin{tabular}{@{\hspace*{75pt}}lp{\Params}}%
\hspace*{30pt}&\\[\Up]%
}
{%
\end{tabular}\end{minipage}\par\vspace{5pt}%
}



\newlength{\bboxw}

\newcommand{\Bbox}[2]{\setlength{\bboxw}{#2pt}\addtolength{\bboxw}{\bboxw}%
\addtolength{\bboxw}{#2pt}%
\addtolength{\bboxw}{\bboxw}\addtolength{\bboxw}{\bboxw}\addtolength{\bboxw}{\bboxw}%
\framebox[\bboxw]{{\footnotesize #1\rule[-4pt]{0pt}{14pt}}}%
}

\newenvironment{vwg}{\par\vspace{10pt}%
\noindent\begin{tabular}{llll}\hspace*{80pt}&\hspace*{200pt}\\[\Up]}
{\end{tabular}\par\vspace{10pt}}


\def\registered{{\ooalign{\hfil\raise.07ex\hbox{\sc r}\hfil\crcr\mathhexbox20D}}}

\newcommand{\Pentium}{Pentium$^{\footnotesize\registered}$}

\setlength{\unitlength}{0.5pt}

\newcommand{\repl}{\leftarrow}
\newcommand{\PT}{P}


\renewcommand{\textfraction}{0.05}
\renewcommand{\topfraction}{0.95}
\renewcommand{\bottomfraction}{0.95}

\hyphenation{ca-che}
\hyphenation{bench-mark}

%%%%%%%%%%%%%%%%%%%%%%%%%%%%%%%%%%%%%%%%%%%%%%%%%%%%%%%%%%%%%%%%%%%%%%
\newcommand{\impnote}[1]{\framebox{\textbf{MIPS Implementation Note:}} #1}
\newcommand{\reg}[1]{\mbox{\textbf{#1}}}
\newcommand{\regs}[3]{{\em #2}&{\footnotesize \reg{#1}}&&{\footnotesize
    \reg{#1}}&{\em #3}\\}
\newcommand{\regx}[4]{{\em #2}&{\footnotesize \reg{#1}}&&{\footnotesize
    \reg{#1}}&{\em #3}&~~{\scriptsize /}~~~{\em #4}\\}

\newcommand{\abox}[2]{\setlength{\bboxw}{#2pt}\addtolength{\bboxw}{\bboxw}%
\addtolength{\bboxw}{\bboxw}\addtolength{\bboxw}{\bboxw}%
\framebox[\bboxw]{{\footnotesize #1\rule[-1ex]{0pt}{4ex}}}}

\newcommand{\bbox}[2]{\setlength{\bboxw}{#2pt}\addtolength{\bboxw}{\bboxw}%
\addtolength{\bboxw}{\bboxw}\addtolength{\bboxw}{\bboxw}%
\framebox[\bboxw]{{\footnotesize #1$_{\ (#2)}$\rule[-1ex]{0pt}{4ex}}}}

\newcommand{\cbox}[3]{\setlength{\bboxw}{#3pt}\addtolength{\bboxw}{\bboxw}%
\addtolength{\bboxw}{\bboxw}\addtolength{\bboxw}{\bboxw}%
\framebox[\bboxw]{{\footnotesize #1$_{\ (#2)}$\rule[-1ex]{0pt}{4ex}}}}

\newenvironment{SC}[1]%
{%
\ifintel\clearpage\markboth{\uppercase{#1}}{\uppercase{#1}}\addcontentsline{toc}{subsection}{#1}\fi%
\vspace*{30pt}\noindent%
\ifintel%
%
{\LARGE\bf #1\\[\Up]}%
\noindent\hspace*{\fill}\begin{tabular}{rl|c|ll@{}l}%
\hspace*{\scindent}&&&&\multicolumn{2}{c}{\hspace*{120pt}}\\&&&&\\&&&&\\%
&&{\Large $-$} \reg{AT} 0x\syscode {\Large $\rightarrow$}&&\\[\Upppp]%
\else%
{\Large\bf #1\\[\Up]}%
\noindent\hspace*{\scindent}\begin{tabular}{l}%
\fi%
}%
{%
\end{tabular}\par\vspace{20pt}%
}

\newenvironment{SC*}[1]%
{%
\vspace*{20pt}%
\noindent\begin{minipage}{\textwidth}%
\noindent{\large\sc #1\\}%
\ifintel%
%
\noindent\hspace*{\fill}\begin{tabular}{rl|c|ll}%
\hspace*{180pt}&&&&\hspace*{120pt}\\&&&&\\&&&&\\%
&&{\large $-$} \reg{AT} 0x\syscode {\Large $\rightarrow$}&&\\[\Upppp]%
\else%
\noindent\hspace*{\scindent}\begin{tabular}{l}%
\fi%
}%
{%
\end{tabular}\end{minipage}\par\vspace{20pt}%
}

%%%%%%%%%%%%%%%%%%%%%%%%%%%%%%%%%%%%%%%%%%%%%%%%%%%%%%%%%%%%%%%%%%%%%%

\pagestyle{headings}

\pagenumbering{roman}
\begin{document}

\thispagestyle{empty}
\setcounter{page}{1}


\maketitle

\setcounter{page}{2}
\pagestyle{headings}
\thispagestyle{empty}

\cleardoublepage

\begin{abstract}
\thispagestyle{headings} \setcounter{page}{3} This document is a user
manual for the L4 \micro-kernel. It gives an introduction to the main
concepts and features of L4, and explains their use by a number of
examples. The manual is generally platform independent, however,
examples are based on the C interface for L4/MIPS. Actual system call C
bindings and data formats differ slightly on other platforms.

This document supplements, rather than replaces, the L4 Reference
Manual, and anyone intending to write an applications on top of L4
should obtain the L4 Reference Manual for their particular platform.

\vspace*{\fill}
{\small
\noindent\hrulefill\\
%
Permission to make digital/hard copy of part or all of this work for
personal or classroom use is granted without fee provided that copies
are not made or distributed for profit or commercial advantage, the
copyright notice, the title of the publication and its date appear, and
notice is given that copying is by permission of the authors. To copy
otherwise, to republish, to post on servers, or to redistribute to lists
requires prior specific permission and/or a fee.\\[2ex]
%
Copyright \copyright 1998 by Gernot Heiser, The University of New South Wales.
}
\end{abstract}
\clearpage
\setcounter{page}{4}
\cleardoublepage

\setcounter{page}{5}
\section*{Preface}

This manual grew out of lecture notes for a course on Advanced Operating
Systems given by Gernot Heiser at UNSW in July--November 1997. In that
course students had to build a small operating system on top of L4/MIPS
which had previously been developed at UNSW by Kevin Elphinstone with
lots of help from Jochen Liedtke. We would like to thank the 44 students
of that course (one of whom was Alan Au) for their interest and
patience, as well as for their questions which prompted the lecturer to
write and provide further documentation. We would further like to thank
Jochen Liedtke for providing the L4 specification and the original
reference manual, Kevin Elphinstone for the MIPS implementation as well
as for explaining L4's idiosyncrasies, the OS group at the Technical
University of Dresden under Herrman H\"artig for example code, C
bindings and manual pages, and Jerry Vochteloo and many others for their
contributions.

It is expected that this manual will continue to grow as more people are
learning to use L4, and consequently discover shortcomings of the
documentation.  We welcome comments, corrections and enhancements and
will try to incorporate them quickly into the manual.  The latest
version of this manual (as well as of the L4/MIPS Reference
Manual~\cite{Elphinstone_HL:L4rm}) are available from
\textsf{http://www.cse.unsw.edu.au/{\Tilde}disy}.

\cleardoublepage
\setcounter{page}{7}
\tableofcontents

\cleardoublepage
\setcounter{page}{1}
\pagenumbering{arabic}
% Permission to make digital/hard copy of part or all of this work for
% personal or classroom use is granted without fee provided that copies
% are not made or distributed for profit or commercial advantage, the
% copyright notice, the title of the publication and its date appear, and
% notice is given that copying is by permission of the authors. To copy
% otherwise, to republish, to post on servers, or to redistribute to lists
% requires prior specific permission and/or a fee.
%
% Copyright (C) 1998 by Gernot Heiser, The University of New South Wales.


% L4 user manual introductory chapter

\chapter{Introduction}

L4 is an operating system microkernel (\micro-kernel). That is, L4 by
itself is not an operating system (OS), but rather constitutes a minimal
base on which a variety of complete operating systems can be built. In
this chapter we will summarise the main ideas behind L4.

The basic idea of a \micro-kernel goes back to Brinch Hansen's
Nucleus~\cite{Brinch_Hansen_70} and Hydra~\cite{Wulf_CCJLPP_74} and has
been popularised by Mach~\cite{Rashid_TYGBBBC_88}. The argument goes
that by reducing the size of the kernel, the part of the operating
system (OS) executing in privileged mode, it becomes possible to build a
system which is more secure and reliable (because the \emph{trusted
computing base} is smaller) and easy to extend. A further advantage is
that a
\micro-kernel-based system can easily implement a number of different
APIs (also called \emph{OS personalities}) without having to emulate one
within the other.

There was also hope of improved efficiency, as operating systems tend to
grow as new features are added, resulting in an increase of the number
of layers of software that need to be traversed when asking for
service. (An example is the addition of the VFS layer in UNIX for
supporting NFS.) A microkernel based system, in contrast, would grow
horizontally rather than vertically: Adding new services means adding
additional servers, without lengthening the critical path of the most
frequently used operations.

However, performance of these first-generation microkernels proved
disappointing, with applications generally experiencing a significant
slowdown compared to a traditional (``monolithic'') operating
system~\cite{Chen_Bershad_93}. Liedtke, however, has
shown~\cite{Liedtke_93, Liedtke_95, Liedtke_96} that these performance
problems are not inherent in the microkernel concept and can be overcome
by good design and implementation. L4 is the constructive proof of this
theorem, as has been clearly demonstrated by H\"artig \emph{et
al.}~\cite{Hartig_HLSW_97}.


\section{L4 design philosophy}

The most fundamental task of an operating system is to provide secure
sharing of resources, in essence this is the only reason why there
\emph{needs to be} an operating system. A \micro-kernel is to be as small
as possible. Hence, the main design criterion of the \micro-kernel is
\emph{minimality} with respect to security: \emph{A service (feature) is
to be included in the \micro-kernel if and only if it impossible to
provide that service outside the kernel without loss of security.} The
idea is that once we make things small (and do it well), performance
will look after itself.

A strict application of this rule has some surprising consequences. For
example device drivers: Some device drivers access physical memory
(e.g.\ DMA) and can therefore break security. They need to be
trusted. This, however, does not mean that they need to be in the
kernel. If they need not execute privileged instructions, and if the
kernel can provide sufficient protection to run them  at user level,
then this is what should be done, and, consequently, this is what L4
does.



\section{L4 abstractions and mechanisms}

Based on such reasoning Liedtke concludes that the \micro-kernel needs
to provide:
\begin{description}
\item[address spaces] because they are the basis of protection,
\item[threads] because there needs to be an abstraction of program
execution,
\item[inter-process communication] (IPC) as there needs to be a way to
transfer data between address spaces,
\item[unique identifiers] (UIDs) for context-free addressing in IPC operations.
\end{description}
We will look at these in turn.


\subsection{Address spaces}

An address-space contains all the data (other than hardware registers)
which are directly accessible by a thread. An address space is a set of
mappings from virtual to physical memory (which is \emph{partial} in the
sense that many mappings are undefined, making the corresponding virtual
memory inaccessible). Address spaces in L4 can be recursively
constructed: A thread can map parts of its address space into another
thread's address space (provided the receiver cooperates) and thereby
share the data mapped by that region of the address space. Mappings can
be revoked at any time, so the mapper retains full control.

Alternatively, virtual address space can be \emph{granted} to a
different address space. In this case the granter relinquishes all
control over the data mapped by that part of the address space and no
longer has a valid mapping for that address space region. The granter
cannot revoke a grant. The grantee, in contrast, inherits full control
over the granted address space (unless the grant was read-only, in which
case write access is lost.) Note that while a grant is irreversible, the
granter has, in general, received the address space (directly or
indirectly) via mapping, and an address space at the beginning of the
mapping chain can still revoke the mapping.

Mapping and granting are implemented as operations on page tables,
without copying any actual data. Mapping and granting is achieved as a
side effect of IPC operations and specified by the means of \emph{flex
pages}. This is not accidental: For security reasons mapping requires an
agreement between sender and receiver, and thus requires IPC
anyway. Details will be explained in Section~\ref{s:fpages}.

The concept of a \emph{task} is essentially equivalent to that of an
address space. In L4, a task is the set of threads sharing an address
space. Creating a new task creates an address space with one running
thread.

Strictly speaking the number of tasks is a constant. There are two kinds
of tasks: \emph{active} and \emph{inactive} ones. When we say that a
task is created we mean that an inactive task is activated. Inactive
tasks are essentially capabilities (task creation rights). This is
important, as a thread can only create a task if it already owns the
task ID to use. Inactive tasks can be donated to other tasks.



There is a hierarchy of tasks, with parents having some limited control
over their children. The main purpose of this is to be able to control
(IPC-based) information flow between address spaces. It has nothing to
do with process hierarchies a particular L4-based personality OS may
implement. Such a hierarchy is under full control of the particular OS
personality.


\subsection{Threads}

A thread is the basic execution abstraction. A thread has an address
space (shared with the other threads belonging to the same task), a UID,
a register set (including an instruction pointer and a stack pointer), a
page fault handler (pager), and an exception handler. IPC operations are
addressed to threads (via their UIDs). Threads are extremely
light-weight and cheap to create, destroy, start and stop. The
lightweight thread concept, together with very fast IPC, is the key to
the efficiency of L4 and OS personalities built on top.


\subsection{IPC}

Message-passing IPC is the heart of L4. The \micro-kernel provides a
total of seven system calls (IPC being one of them), which provide some
very rudimentary OS functionality. Everything else must be built on top,
implemented by server threads, which communicate with their clients via
IPC.

IPC is used to pass data by value (i.e., the \micro-kernel copies the
data between two address spaces) or by reference (using mapping or
granting). IPC is also used for synchronisation (as it is blocking, so
each successful IPC operation results in a \emph{rendez-vouz}),
wakeup-calls (as timeouts can be specified), pager invocation (the
\micro-kernel converts a page fault into an IPC to a user-level pager),
exception handling (the \micro-kernel converts an exception fault into
an IPC to a user-level exception handler), and interrupt handling (the
\micro-kernel converts an interrupt fault into an IPC to a user-level
interrupt-handler from a pseudo-thread). Device control is registered
via IPC (although actual device access is memory mapped).


\subsection{\label{s:i:c&c}Clans \& chiefs}

Clans and chiefs are L4's basic mechanism enabling the implementation of
arbitrary security policies~\cite{Liedtke_92}. They allow controlling
IPC and thus information flow.

The basic idea is simple: A task's creator is that task's \emph{chief},
all tasks (directly) created by a particular chief constitute that
chief's \emph{clan}. Threads can directly send IPC only to other threads
in the same clan, or to their chief. If a message is sent to a thread
outside the clan containing the sender, that message is instead
delivered to the sender's chief (who may or may not forward the
message). If a message is sent to a member of a subclan of the clan
containing the sender, that message is delivered to the task in the clan
whose clan (directly or indirectly) contains the addressee.

This is depicted in Fig.~\ref{f:i:c&c}, where circles represent tasks,
ovals clans (with their chief on top), and arrows symbolise IPC. The
bold arrow indicates the \emph{intended} IPC, while the thin arrows
indicate the way the IPC is actually routed (provided all chiefs
cooperate).

\begin{figure}[htb]

\setlength{\unitlength}{0.5pt} 
\begin{center}\begin{picture}(320,240)(0,-50) 
 
\put(170,170){\circle{20}} 
\put(70,120){\circle{20}} 
\put(90,60){\circle{20}} 
\put(40,30){\circle{20}} 
 
\put(250,120){\circle{20}} 
\put(200,85){\circle{20}} 
\put(200,25){\circle{20}} 
\put(270,40){\circle{20}} 
 
\thicklines 
\put(105,55){\vector(3,-1){80}} 
\thinlines
 
\put(85,75){\vector(-1,3){10}} 
\put(85,120){\vector(1,0){150}} 
\put(245,110){\vector(-2,-1){30}} 
\put(200,70){\vector(0,-1){30}} 
 
\put(70,60){\oval(100,100)} 
 
\put(200,40){\oval(60,70)} 
 
\put(250,50){\oval(190,120)} 

\put(170,60){\oval(380,200)} 

\end{picture}\end{center} 
\caption{\label{f:i:c&c}Message redirection by clans \& chiefs}
\end{figure}


\subsection{UIDs}

A UID of a thread is that of its task plus the number of the thread
within the task (called \emph{lthread} or local thread number). The UID
of a task consists of the task number, some fields describing its place
in the task hierarchy, and a version number.

Both, tasks and threads are limited (on L4/MIPS there are 2048 tasks and
within each task 128 threads). This means that tasks (or address spaces)
and threads must be recycled. The \micro-kernel ensures uniqueness of
task IDs by incrementing the version number whenever a task number is
reused.

As far as threads are concerned, the L4 view is that threads do not die
as long as their task exists, they can only be blocked (waiting for IPC
which will never arrive). This avoids the issue of lthread uniqueness.

Obviously, both thread and task numbers are insufficient for a real
multi-user operating system. This means that an OS personality will in
general need to map its own task and thread abstraction onto L4's. How
this is done is up to the OS personality, L4 only provides the tools.



\section{Resource allocation}

As said earlier, the classical job of the OS is resource allocation. In
a \micro-kernel-based system, this is left to the OS personality, the
\micro-kernel only provides the tools, and enforces security.

L4's view of resources is simple: Each resource is allocated on a
first-come-first-served basis. This is by no means a free-for-all: The
servers implementing OS personalities are started at boot time by the
\micro-kernel. As they are the first running tasks, they have the chance
to allocate all resources to themselves before any ``user'' tasks
exist. (As most resources are claimed via IPC the clans \& chiefs
mechanism would anyway prevent direct access to resources by tasks which
are not top-level.)

Initial servers must be contained in the L4 boot image, and must be
identified as to be started up by L4. This leaves responsibility for
providing ``sensible'' servers to whoever creates the boot
image. Obviously, if you create the boot image you are in full control
over what gets to run, and what is in the system. Hence this approach is
secure.

The same argument holds for ``competition'' between initial servers (OS
personalities). This is under full control of the system designer (and
whoever creates the boot image) and therefore secure.


\section{Organisation of this manual}

Having presented the basic ideas on which L4 is based the next chapters
will discuss L4 in detail. Chapter~\ref{c:ipc} explains the IPC system
call, which is by far the most complex one in L4. The remaining L4
system calls are covered in Chapter~\ref{c:other}. Chapter~\ref{c:use}
discusses general use of the system.

% 
%%% Local Variables: %%%
%%% mode: latex %%%
%%% End: %%%


% Permission to make digital/hard copy of part or all of this work for
% personal or classroom use is granted without fee provided that copies
% are not made or distributed for profit or commercial advantage, the
% copyright notice, the title of the publication and its date appear, and
% notice is given that copying is by permission of the authors. To copy
% otherwise, to republish, to post on servers, or to redistribute to lists
% requires prior specific permission and/or a fee.
%
% Copyright (C) 1998 by Gernot Heiser, The University of New South Wales.


% L4 user manual IPC chapter
% Author: Alan Au
% Created On: 3 December 1997
% Last Modified: 3 Septermber 1998

\chapter{\label{c:ipc}L4 IPC}


%%%%%%%%%%%%%%%%%%%%%%%%%%
%%%%%%%%%%%%%%%%%%%%%%%%%%

\section{IPC Overview}

%%%%%%%%%%%%%%%%%%%%%%%%%%
%%%%%%%%%%%%%%%%%%%%%%%%%%

Message passing is the basic IPC mechanism in L4. It allows L4 threads
to communicate via messages.

All L4 IPC is synchronous and unbuffered. Synchronous IPC requires an
agreement between both the sender and the receiver. The main
implications of this agreement is that the receiver is expecting an
IPC and provides the necessary buffers. If either the sender or the
receiver is not ready, then the other party must wait. Unbuffered IPC
reduces the amount of copying involved and is thus the key to high
performance IPC.

%%%%%%%%%%%%%%%%%%%%%%%%%%
%%%%%%%%%%%%%%%%%%%%%%%%%%

\section{L4 IPC Messages}

%%%%%%%%%%%%%%%%%%%%%%%%%%
%%%%%%%%%%%%%%%%%%%%%%%%%%

%%%%%%%%%%%%%%%%%%%%%%%%%%
%%%%%%%%%%%%%%%%%%%%%%%%%%

\subsection{Message Data Types}

%%%%%%%%%%%%%%%%%%%%%%%%%%
%%%%%%%%%%%%%%%%%%%%%%%%%%

\label{sec:msgdata}

Data can be transferred in three ways using L4 IPC.

\begin{enumerate}

\item In-line by-value data. A limited amount of such data is passed
  directly in registers (first 8 words in MIPS R4k) with any remainder
  in a message buffer.

\item Strings. Arbitrary out-of-line buffers which are copied to the
  receiver.
        
\item Virtual memory mappings (by-reference data). Data transfer via
  mappings is described by \emph{flex-pages (fpages)}. Alternatively,
  virtual memory can be granted: mapped to the receiver and unmapped
  from the sender simultaneously.

\end{enumerate}

Figure~\ref{fig:msgdata} illustrates the difference between in-line
by-value data and strings.

\begin{figure}[htb]
\centering
\mbox{\subfigure[In-Line]
        {\epsfig{figure=eps/inline_data.eps,width=.32\textwidth}}\quad
      \subfigure[String]
        {\epsfig{figure=eps/string_data.eps,width=.22\textwidth}}}
\caption{In-line Versus String Data}
\label{fig:msgdata}
\end{figure}

A \label{s:fpages}flex-page is a contiguous region of virtual address
space. A flex-page has the following properties.

\begin{itemize}
        
\item Of size \(2^{s}\) bytes. The smallest size allowed for any fpage
  is the hardware page size.

\item Base address aligned to \(2^{s}\) (\(base~address \bmod~2^{s} =
  0 \)).

\item Contains all the \emph{mapped} pages within the region described
  by the flex-page. That is, those pages which belong to the specified
  region \emph{and} which have been mapped into the sender's virtual
  address space. (For convenience, call these the valid pages).

\end{itemize}

Fpages are required for mapping and granting virtual memory. Fpages
are specified by the mapper and received by the mappee\footnote{The
  terms \emph{mapper} and \emph{mappee} will be used in both mapping
  and granting contexts.} as part of an IPC message. For each fpage
successfully received, the valid pages within that fpage become part
of (mapped or granted to) the receiver's address space.

\begin{figure}[htb]
\begin{center}
\leavevmode
\epsfxsize = .7\textwidth
\epsfbox{eps/fpage.eps}
\end{center}
\caption{Flex-page}
\label{fig:fpage}
\end{figure}

An fpage is specified by providing the values $b$ and $s$. The fpage
is then defined to be the region [$b \times 2^{s}$, $(b + 1) \times
2^{s}$]. In addition, a hot-spot, $h$, is also required for the sender
if the sender and receiver specify fpages of different sizes. In such
a case, the hot-spot specification is used to determine how the
mapping between the two different size fpages occurs: If \(2^s\) is
the size of the larger, and \(2^t\) the size of the smaller fpage,
then the larger fpage can be thought of as being tiled by \(2^{s-t}\)
fpages of the smaller size. One of these is uniquely identified as
containing the hot spot address (mod \(2^s\)). This is the fpage which
will actually be mapped/granted.

Figure~\ref{fig:fpmap} gives two examples of mappings which involve
different fpage sizes. The figure also illustrates the fact that the
receiver's fpage allows the receiver to specify the window where
mappings are permitted to occur (greater security).

\begin{figure}[htb]
\begin{center}
\leavevmode
\epsfxsize = .7\textwidth
\epsfbox{eps/fpmap.eps}
\end{center}
\caption{Fpage Mapping Example}
\label{fig:fpmap}
\end{figure}

A precise definition of the intuitive description above is now given
for completeness. If the sender defines its fpage as $b$, $s$, $h$ and
the receiver defines its fpage as $b'$, $s'$ then the mappings for the
three situations would be:

\begin{description}

\item[$s = s'$:] mapping is $b \times 2^{s} \mapsto b' \times 2^{s}$
  \\ hot spot is not used

\item[$s < s'$:] mapping is $b \times 2^{s} \mapsto b'_{[63, s]}
  h_{[s' - 1, s]}0_{(s)}$ \\ the sender's fpage is aligned around the
  hot-spot

\item[$s > s'$:] mapping is $b_{[63, s]}h_{[s - 1, s']}0_{(s')}
  \mapsto b' \times 2^{s'}$\\ the receiver's fpage is aligned around
  the hot-spot

\end{description}

As an explanation of the notation above, $b'_{[63, s']}h_{[s' - 1,
  s]}0_{(s)}$ represents the bit address formed by concatenating
together, in the given order:

\begin{itemize}

\item Bits 63 through to bit $s'$ (inclusive) of $b$.
        
\item Bits $s' - 1$ through to bit $s$ (inclusive) of $h$.

\item $s$ number of zeroes.

\end{itemize}   

%%%%%%%%%%%%%%%%%%%%%%%%%%
%%%%%%%%%%%%%%%%%%%%%%%%%%

\subsection{Messages}

%%%%%%%%%%%%%%%%%%%%%%%%%%
%%%%%%%%%%%%%%%%%%%%%%%%%%

\label{sec:messages}

The message formats shown in this section are for L4/MIPS. On other
platforms they are mostly very similar.

%%%%%%%%%%%%%%%%%%%%%%%%%%
%%%%%%%%%%%%%%%%%%%%%%%%%%

\subsubsection{Message Descriptors}

%%%%%%%%%%%%%%%%%%%%%%%%%%
%%%%%%%%%%%%%%%%%%%%%%%%%%

A message descriptor is a pointer to the start of a message
buffer\footnote{See following description of a long message.} or
indication that the IPC is purely register based. There are two types
of message descriptors: one for sending and one for receiving IPC.

The format of the send message descriptor is:

\begin{tabular}{ll} 
\hspace*{100pt}&\\[\Up] 
% 
&
\cbox{snd msg/4}{62}{30}\abox{m}{1}\abox{d}{1}\\[20pt]
% 
\end{tabular}

A non-zero message descriptor address (\emph{snd msg}) is interpreted
as the start address of the sender's message buffer. A zero value
indicates a purely register based IPC. Setting the $m$-bit indicates
that the IPC includes mappings (i.e.\ fpages are present followed
possibly by by-value data). A zero value for the $m$-bit indicates
that the message contains only by-value data and no fpages. Setting
the $d$-bit indicates that the IPC is deceiving.

The format of the receive descriptor is similar:

\begin{tabular}{ll} 
\hspace*{100pt}&\\[\Up] 
% 
&
\cbox{rcv msg/4}{62}{30}\abox{m}{1}\abox{o}{1}\\[20pt]
% 
\end{tabular}

If the $m$-bit is not set, the message descriptor address (\emph{rcv
  msg}) is interpreted as the start address of the receiver's message
buffer, which may contain a receive fpage, indicating the caller's
willingness to receive mappings or grants. Setting the $m$-bit
indicates that the caller is willing to accept fpage mappings but no
long message, and has supplied the receive fpage directly as the
\emph{rcv msg} parameter (there is no message buffer in this case).
Setting the $o$-bit will allow a receive from any sender (open wait)
while a zero value for the $o$-bit allows receiving only from the
specified sender.

Note that the the message descriptor addresses have had their least
significant two bits removed. These two bits are not needed as the
message buffer must be word aligned.

%%%%%%%%%%%%%%%%%%%%%%%%%%
%%%%%%%%%%%%%%%%%%%%%%%%%%

\subsubsection{Short And Long Messages}

%%%%%%%%%%%%%%%%%%%%%%%%%%
%%%%%%%%%%%%%%%%%%%%%%%%%%

Every successful IPC operation will always copy at least eight dwords
(for MIPS) to the receiver. These eight dwords contain the first 64
bytes of a message's in-line data\footnote{Section~\ref{sec:msgdata}
  describes the two types of in-line data: by-value data and fpage
  descriptors.} and is referred to as the \emph{short} part of the
message. The short message is transferred via registers (\textbf{s0}
\dots \textbf{s7} on MIPS R4k) and its format is:

\vspace{15pt}
\noindent\begin{tabular}{lcr} 
\hspace*{100pt}
s7 & \abox{by-value data}{32} & \\
\hspace*{100pt} \vdots & \vdots & \\
\hspace*{100pt} s0 & \abox{fpage descriptors}{32} &
\end{tabular}

The presence of fpages is indicated by setting the $m$-bit in the
message descriptor. Processing of fpages starts at the beginning of
the message and continues until an invalid fpage is encountered. This
last dword and the remainder of the in-line data is interpreted as
by-value data.

The \emph{long} part of the message is optional and its presence is
indicated by the message descriptor (\emph{snd msg/recv msg}). If
present, it is a dword-aligned memory buffer pointed to by a
\emph{message descriptor}. The buffer contains a three dword message
header, followed by a number of mwords (the rest of the in-line data),
followed by a number of string dopes. The number of mwords (in 64-bit
dwords, excluding those copied in registers) and string dopes is
specified in the message header.

The format for the long part of the message is:

 \vspace{15pt}
\noindent\begin{tabular}{lcr} 
 \hspace*{100pt}  & \abox{string dopes}{32} & +$x$\\
                  & \abox{mwords (fpage descriptors + by-value data)}{32} & +24\\
message:          & \abox{msg header}{32} & +0\\
\end{tabular}

The value of $x$ is determined by the number of mwords in the message
as specified in the \emph{size dope} of the message header.

%%%%%%%%%%%%%%%%%%%%%%%%%%
%%%%%%%%%%%%%%%%%%%%%%%%%%

\subsubsection{Message Header}

%%%%%%%%%%%%%%%%%%%%%%%%%%
%%%%%%%%%%%%%%%%%%%%%%%%%%

The message header describes the format of the long message.

 \vspace{15pt}
\noindent\begin{tabular}{lcr} 
msg snd dope:  &
          \cbox{0}{32}{9}\cbox{dwords}{19}{10}\cbox{strings}{5}{5}\cbox{\undef}{8}{8} & +16\\
msg size dope:  &
          \cbox{0}{32}{9}\cbox{dwords}{19}{10}\cbox{strings}{5}{5}\cbox{\undef}{8}{8} & +8\\
msg rcv fpage option: & \cbox{fpage}{64}{32} & +0\\
\end{tabular}

The \emph{size dope} defines the size of the dword buffer, in words,
(and hence the offset of the string dopes from the end of the header),
and the number of string dopes in the long message.

The \emph{send dope} specifies how many dwords and strings are
actually to be sent. (Specifying send dope values less than the size
dope values makes sense when the caller is willing to receive more
data than it is sending.)

The \emph{receive fpage} describes the address range in which the
caller is willing to accept fpage mappings or grants in the receive
part (if any) of the IPC. As described in section~\ref{sec:msgdata},
an fpage region is defined by providing its base address, $b$ and size
exponent, $s$. Note that the hot-spot, $h$, is provided by the sender
and hence not required as part of the receive fpage.

The fpage format is:  

\vspace{12pt} 
 
\begin{tabular}{ll} 
\hspace*{80pt}&\\[\Up] 
% 
fpage$(b,2^{s})$ &
\cbox{$b/4096$}{52}{20}\bbox{0}{3}\cbox{$s$}{7}{7}\abox{\undef{}}{1}\abox{\undef{}}{1}\\[20pt]
% 
\end{tabular} 

The base address can be given as $b/4096$ rather than just $b$ because
fpages needs to be aligned to at least the hardware page size ($4096$
or $2^{12}$ bytes).

A special case of an fpage is one that specifies the complete user
address space (i.e. address space with base 0 and size $2^{64} - K$,
where $K$ is the size of the kernel area).

\vspace{12pt} 
 
\begin{tabular}{ll} 
\hspace*{80pt}&\\[\Up] 
% 
fpage$(0,2^{64} - K)$ &
\cbox{0}{52}{20}\bbox{0}{3}\cbox{64}{7}{7}\abox{\undef{}}{1}\abox{\undef{}}{1}\\[20pt]
% 
\end{tabular} 

%%%%%%%%%%%%%%%%%%%%%%%%%%
%%%%%%%%%%%%%%%%%%%%%%%%%%

\subsubsection{Message mwords}

%%%%%%%%%%%%%%%%%%%%%%%%%%
%%%%%%%%%%%%%%%%%%%%%%%%%%

The (possibly zero) message mwords follow directly after the message
header and contain the rest of the in-line data remaining after the
short message. This in-line data is made up of a number (again,
possibly zero) of fpage descriptors followed by by-value data.

 \vspace{15pt}
\noindent\begin{tabular}{lcr} 
\hspace*{100pt} & \vdots & \\
& \abox{by-value data}{32} & \\
mwords:   & \abox{fpage descriptors}{32} &\\
& \vdots &
\end{tabular}

Fpage descriptors are expected in memory (the long message) only if
the $m$-bit is set in the message descriptor and all register data
(the entire short message) consists of valid fpages. These fpage
descriptors (together with those in the short message) are provided by
the sender for memory mapping purposes. The format of an fpage
descriptor is:

\vspace{12pt} 

\label{s:w-bit}
\begin{tabular}{ll} 
\hspace*{80pt}&\\[\Up] 
% 
snd fpage: &
\cbox{$b/4096$}{52}{20}\bbox{0}{3}\cbox{$s$}{7}{7}\abox{w}{1}\abox{g}{1}\\
%
hot spot: &\cbox{snd base}{64}{32} 
\end{tabular} 

The \emph{snd fpage} is in the same format as the \emph{receive fpage}
in the message header except the least significant two bits are no
loner undefined. Setting the $w$ bit will cause the fpage to be mapped
writable (rather than just read-only) and setting the $g$ bit causes
the fpage to be granted (rather than just mapped). The \emph{snd base}
is the mapping hot spot as described in section~\ref{sec:msgdata}.

The kernel will interpret each pair of dwords of the in-line part
(starting from the short message part and continuing to the long
message part, if present) as fpage descriptors until an invalid
descriptor is encountered. This and any further dwords are then
interpreted as by value data.

Note that all in-line data is copied to the receiver, including any
initial parts which are interpreted as fpage descriptors.

%%%%%%%%%%%%%%%%%%%%%%%%%%
%%%%%%%%%%%%%%%%%%%%%%%%%%

\subsubsection{String Dopes}

%%%%%%%%%%%%%%%%%%%%%%%%%%
%%%%%%%%%%%%%%%%%%%%%%%%%%

The last component making up a long message are string dopes. There
can be zero or more string dopes, with the exact number specified in
the \emph{size dope} of the message header. Each string dope describes
a region in memory where out-of-line data can be copied from (on an
IPC send) and copied to (on an IPC receive).  The size and location of
each string is specified in a string dope. The kernel copies data from
sender memory, specified by the sender string dopes, to the receiver's
memory, specified by the receiver's string dopes. Each string dope
occupies four dwords and its format is:

\vspace{15pt}
\noindent\begin{tabular}{lcr} 
\hspace*{100pt}& \cbox{*rcv string}{64}{32} & \\
& \cbox{rcv string size}{64}{32} &\\
& \cbox{*snd string}{64}{32} & \\
& \cbox{snd string size}{64}{32} &
\end{tabular}

The first part of the string dope specifies the size and location of
the string the caller wants sent to the destination, while the second
part specifies the size and location of a buffer where the caller is
willing to receive a string. \textbf{Note} that strings do not have to
be aligned, and that their size is specified in \emph{bytes}.


%%%%%%%%%%%%%%%%%%%%%%%%%%
%%%%%%%%%%%%%%%%%%%%%%%%%%

\subsection{Sending/Receiving IPC Messages}

%%%%%%%%%%%%%%%%%%%%%%%%%%
%%%%%%%%%%%%%%%%%%%%%%%%%%

Before considering \emph{how} to send or receive it is necessary to
decide \emph{what} form of data is to be sent. In particular, a
decision must be made on whether to send data in-line or as strings.
Both have the same effect of making a copy of the sender's data
available to the receiver. The difference lies in efficiency and
appropriateness.

In-line data needs to be copied to a buffer first and must also be
aligned. Thus the in-line option is best for small amounts of data and
is useful when some marshaling is required anyway. Sending short
strings as in-line data is more efficient as it avoids setting up
string dopes and may also be done in registers.

Strings avoid extra copying and can be located anywhere in memory (no
alignment necessary) but require buffers to be specified (via string
dopes). Buffer specifications must also be consistent on both the
sender and the receiver end. In particular, the receiver must specify
and expect to receive the (at least) same number and size strings that
the sender sends.

The C interface to L4's system calls provides a number of IPC
operations. They are differentiated in the following ways:

\begin{itemize}

\item Single send or receive versus a combined send and receive.
        
\item Receive from specific sender (closed) or any sender (open).

\item Deceiving or non-deceiving IPC.

\end{itemize} 

To send or receive a message, certain parameters must be provided.

\begin{itemize}
        
\item \emph{dest/src} Identifier of message destination/source thread
  respectively.
        
\item \emph{snd\_msg/rcv\_msg} Descriptor for long part of message,
  for send/receive part of IPC respectively.

\item \emph{snd\_reg/rcv\_reg} Short part of message (artifact of C
  interface) for send/receive part of IPC respectively.

\item \emph{timeout} Timeout specification. Used to ensure that a
  thread need not be blocked indefinitely.

\item \emph{result} Variable to store the result status of the IPC
  operation.

\end{itemize}


%%%%%%%%%%%%%%%%%%%%%%%%%%
%%%%%%%%%%%%%%%%%%%%%%%%%%

\cbstart
\subsubsection{L4 IPC Message Summary}

In summary, an L4 IPC operation can have a send and a receive
phase.  The \emph{snd\_msg} descriptor describes what is to be sent, and the
\emph{rcv\_msg} descriptor describes what can be received. For the IPC to be
successful, the sender's send descriptor and the receiver's receive
descriptor must be compatible.

Every successful IPC transfers some by-value data, the register (or
short) part of the direct string (8 bytes on ix86, 64 bytes on
R4x00). In addition, the following may be transferred, provided the
\emph{snd\_msg} descriptor says so, and the \emph{rcv\_msg} descriptor
allows it:

\begin{itemize}

\item a bigger direct string, provided that:
\begin{itemize}
\item the send descriptor points to a  message descriptor  specifying a
     non-zero number of dwords in the \emph{send dope}, and
\item the receive descriptor points to a message descriptor specifying a
     non-zero number of dwords in the \emph{size dope}, and
\item there is no error;
\end{itemize}

\item one or more indirect strings, provided that:
\begin{itemize}
\item the send descriptor points to a  message descriptor specifying a
     non-zero number of strings in the \emph{send dope}, and
\item the receive descriptor points to a message descriptor specifying a
     non-zero number of strings in the \emph{size dope}, and
\item there is no error;
\end{itemize}

\item one or more page mappings or grants, provided that
\begin{itemize}
\item the \emph{m}-bit is set in the send descriptor, and
\item the beginning of the sender's direct string (starting with the
     register part) contains at least one valid fpage descriptor, and
\item the receive descriptor either has the form of a valid \emph{receive
     fpage} and has the \emph{m}-bit set, or points to a message
     descriptor containing a valid \emph{receive fpage}, and
\item there is no error.
\end{itemize}
\end{itemize}

Note that the structure of the message descriptor and the string dopes
make it easy to use the same message header for sending and receiving
(i.e., having the \emph{send descriptor} and \emph{receive descriptor}
point to the same address). The \emph{send dope} specifies the size of
the direct string and the number of indirect strings for the send part
(if any) of the IPC, while the \emph{size dope} specifies the size of
the buffer for the direct string and the number of buffers for indirect
strings for the receive operation (if any). Note that using the same
message descriptor for sending and receiving implies using the same
buffer for the direct string. For indirect strings, each string dope
specifies separately the location and size of the buffers for the
strings to be sent and received. If the IPC does not contain a send
\textbf{and} a receive part, then some of the information in the message
header is not used. Similarly, if the send and receive descriptors point
to different message structures, some of the information in them is unused.

Obviously, if the same message descriptor is used for sending and
receiving, receiving a direct string (longer than the register part)
will overwrite the string sent. Similarly, if the \emph{receive string}
of some string dope points to the same address as the \emph{send string}
of the same or another string dope, then receiving may overwrite some of
the data which has been sent. However, as the send part of the IPC is
guaranteed to be concluded before any receive action takes place, this
does not create any problems if the sender does not need the data any
more. The data to be sent will have been safely copied to the receiver
prior to the receive part of the IPC overwriting it on the caller's end.

\cbend

%%%%%%%%%%%%%%%%%%%%%%%%%%
%%%%%%%%%%%%%%%%%%%%%%%%%%

\subsubsection{L4 Timeouts}

%%%%%%%%%%%%%%%%%%%%%%%%%%
%%%%%%%%%%%%%%%%%%%%%%%%%%

Timeouts are used to control IPC operations. Each IPC operation
specifies four timeout values. The first two are with respect to the
time \emph{before message transfer starts}. These timeouts are no
longer relevant once message transfer starts. \\ 
 
\begin{center} 
\bbox{$m_r$}{8}\bbox{$m_s$}{8}\bbox{$p_s$}{4}\bbox{$p_r$}{4}\bbox{$e_s$}{4}\bbox{$e_r$}{4} 
\end{center} 

\begin{description}

\item[receive timeout] Specifies how long to wait for an incoming
  message. The receive operation fails if the timeout period is
  exceeded before message transfer starts. The receive timeout is
  calculated as $m_r4^{15-e_r} \mu s$.

\item[send timeout] Specifies how long IPC should try to send a
  message. The send operation fails if the timeout period is exceeded
  before message transfer starts. The send timeout is calculated as
  $m_s4^{15-e_s} \mu s$.

\end{description}

The other timeout values are used if a page fault occurs during an IPC
operation. A page fault is converted to an RPC to a pager by the
kernel (see section~\ref{sec:pf}). The page fault timeouts are with
respect to this RPC.

\begin{description}

\item[receive page fault timeout] Used for both the send and receive
  timeouts of the page fault RPC when a page fault occurs in the
  sender's address space during an IPC. This value is set by the
  receiver and is calculated as $4^{15-p_r} \mu s$.

\item[send page fault timeout] Used for both the send and receive
  timeouts of the page fault RPC when a page fault occurs in the
  receiver's address space during an IPC. This value is set by the
  sender and is calculated as $4^{15-p_s} \mu s$.

\end{description}

There are two special timeout values: $\infty$ and 0. An infinite
value means no timeout (i.e.\ possibly indefinite blocking) and is
specified by zero values of $e$ or $p$. A zero timeout value
represents non-blocking IPC and is specified by zero values of $m$
(with $e > 0$). A maximum value for $p$ ($p = 15 $) means that the IPC
will fail if a page fault occurs.


%%%%%%%%%%%%%%%%%%%%%%%%%%
%%%%%%%%%%%%%%%%%%%%%%%%%%

\subsubsection{L4 IPC Result Status}

%%%%%%%%%%%%%%%%%%%%%%%%%%
%%%%%%%%%%%%%%%%%%%%%%%%%%

The status of each IPC operation is returned in a message dope with
the following format.\label{s:ccode}

\begin{center}
\cbox{0}{32}{10}\cbox{mwords}{19}{9}\bbox{strings}{5}% 
                          \bbox{$cc$}{8}\\
\end{center}

\begin{description}

\item[mwords] Size \emph{in words} of in-line data received (excluding
  register data).

\item[str] Number of strings received.

\item[cc] Condition code. \\ \\ 
  \bbox{$ec$}{4}\abox{$i$}{1}\abox{$r$}{1}\abox{$m$}{1}\abox{$d$}{1}
  \\ \\ \textbf{ec} is an \emph{error code} associated with the IPC. A
  zero value for \textbf{ec} implies a successful IPC while a non-zero
  value means the IPC failed for some reason. On IPC failure, the
  actual value of \textbf{ec} reports the reason for the failure
  (refer to reference manual). A common cause of IPC failure results
  when the receiver expects less than what is really sent (i.e.\ the
  receive buffer is too small, not providing enough receive strings,
  etc.). Note that the value of \textbf{ec} is also delivered as the
  return value of the IPC procedures in the C library interface.

  In the case of a truncated IPC the \micro-kernel provides no
  information on how much data the sender was trying to transmit.

\end{description}

%%%%%%%%%%%%%%%%%%%%%%%%%%
%%%%%%%%%%%%%%%%%%%%%%%%%%

\subsubsection{Send/Receive Protocol}

%%%%%%%%%%%%%%%%%%%%%%%%%%
%%%%%%%%%%%%%%%%%%%%%%%%%%

In the overview to this chapter, it was mentioned that the sender and
receiver of an IPC must make certain agreements. Sender and receiver
must agree on the following points for the IPC:

\begin{itemize}

\item The size of data to be copied.

\item The number and size of strings to be transferred.

\item Whether the IPC involves memory mappings (the presence of
  fpages).

\end{itemize}

The kernel does not provide the receiver with any information (e.g.
size of data) concerning the incoming message. Thus, a user message
protocol needs to be defined beforehand to ensure agreement on the
above points. In general, the receiver can expect more from the sender
than is actually sent but not less.

Example of such a protocols are:

\begin{itemize}

\item The sender marshals in-line data into a particular structure
  which the receiver will be expecting. Depending on the structure and
  context, the receiver may not make use of the entire structure but
  will receive the entire structure regardless (e.g. one field of the
  structure may identify the context and consequently how the other
  fields are to be interpreted).

\item The receiver always receives a maximum number of maximum size
  strings.

\item Use two IPCs. The first one allows the sender to establish an
  agreement with the receiver. The second IPC is the main message in
  the form established by the first IPC.

\end{itemize} 

%%%%%%%%%%%%%%%%%%%%%%%%%%
%%%%%%%%%%%%%%%%%%%%%%%%%%

\subsubsection{Sending}

%%%%%%%%%%%%%%%%%%%%%%%%%%
%%%%%%%%%%%%%%%%%%%%%%%%%%


The following is a suggested procedure that can be followed to send a message.

\begin{enumerate}

\item Declare a result status variable.

\item Declare a register buffer for the short part of the message. If
  there is any in-line data copy fpages and/or by-value data into the
  buffer.

\item If a long message is required, decide on a message format and
  determine the message descriptor. The format is determined by the
  number of fpages, the amount of by-value data and the number of
  strings to be sent.

  \begin{enumerate}

  \item Fill in the message header.
                
  \item Copy in fpage descriptors (if any).
        
  \item Copy in by-value data (if any).

  \item Copy in (send) string dopes, one for each string.

  \end{enumerate} 

\item Determine the thread id of the desired receiver thread.

\item Determine the desired timeout period.

\item Provide the parameters for one of the C interface procedures
  which allows for an IPC send.

\end{enumerate}

\example{Sending a short message}

This example illustrates sending a short message. Assume the following
conditions:

\begin{itemize}

\item A single non-deceiving IPC send to a specific thread.
\item The receiver thread is $\sigma_0$.
\item Purely register based message (short message).
\item No timeout (indefinite blocking).

\end{itemize}

Then following the suggested procedure:

\begin{enumerate}

\item \emph{Declare a result status variable.} \\ \\
  \hspace*{20pt} {\footnotesize\verb+l4_msgdope_t result;+}

\item \emph{Declare a register buffer.} \\ \\
  \hspace*{20pt} {\footnotesize\verb+l4_ipc_reg_msg_t rmsg;+} \\ \\
  \hspace*{10pt} At this point, in-line data may be copied into the 
  register buffer. For example (Note that the entire buffer need not 
  be used): \\ \\
  \hspace*{20pt} {\footnotesize\verb+rmsg.reg[0] =+} \ldots; {\footnotesize\verb+ /* a dword */+} \\
  \hspace*{20pt} {\footnotesize\verb+rmsg.reg[1] =+} \ldots; {\footnotesize\verb+ /* a dword */+} \\
  \hspace*{50pt} \vdots

\item \emph{If a long message is required, decide on a message format.}
  \\ \\ \hspace*{10pt}
  Not required for a short message. To indicate that the message is 
  purely message based, use {\footnotesize\verb+L4_IPC_SHORT_MSG+} (constant defined in
  {\footnotesize\verb+ipc.h+}) for the message descriptor (\emph{snd\_msg} parameter).

\item \emph{Determine the thread id of the desired receiver thread.} \\
  \\ \hspace*{10pt}
  The thread id of $\sigma_0$ is given by {\footnotesize\verb+SIGMA0_TID+} 
  (constant defined in {\footnotesize\verb+sigma0.h+}).
  
\item \emph{Determine the desired timeout period.} \\ \\ \hspace*{10pt}
  An infinite timeout period is provided by the constant 
  {\footnotesize\verb+L4_IPC_NEVER+} (defined in {\footnotesize\verb+types.h+}). 
  
\item \emph{Provide the parameters for one of the C interface 
    procedures.} \\ \\ \hspace*{10pt}
  For the given conditions the most appropriate C interface procedure 
  is:\\ \\
  \hspace*{20pt} {\footnotesize\verb+int l4_mips_ipc_send(l4_thread_t dest,+} \\ 
  \hspace*{130pt} {\footnotesize\verb+const void *snd_msg,+} \\
  \hspace*{130pt} {\footnotesize\verb+l4_ipc_reg_msg_t *snd_reg,+}\\ 
  \hspace*{130pt} {\footnotesize\verb+l4_timeout_t timeout,+}\\ 
  \hspace*{130pt} {\footnotesize\verb+l4_msgdope_t *result)+}\\ \\
  \hspace*{10pt} Filling in the parameters using information from steps
  1-5 gives the desired system call: \\ \\
  \hspace*{20pt} {\footnotesize\verb+l4_mips_ipc_send(SIGMA0_TID, L4_IPC_SHORT_MSG,+}\\
  \hspace*{140pt} {\footnotesize\verb+&rmsg, L4_IPC_NEVER, &result);+}

\end{enumerate}  

\example{Sending a long message} 

This example illustrates sending a long message that contains in-line
as well as string data. Assume the following conditions:

\begin{itemize}

\item A single non-deceiving IPC send to a specific thread.
\item The receiver thread is $\sigma_0$.
\item Uses in memory message buffer (long message).
\item No fpages.
\item The in-line by-value data to be sent is stored in a variable
  with the following declaration: \\
  \hspace*{20pt} {\footnotesize\verb+dword_t buf[10];+}
\item One string is to be sent. The string is a null terminated
  character string called {\footnotesize\verb+sbuf+}.
\item No timeout (indefinite blocking).
  
\end{itemize}

Again, following the suggested procedure:

\begin{enumerate}

\item \emph{Declare a result status variable.} \\ \\
  \hspace*{20pt} {\footnotesize\verb+l4_msgdope_t result;+}
  
\item \emph{Declare a register buffer and copy data into it.} \\ \\
  \hspace*{10pt} Note that in this case, no direct copying is needed. All
  that is required is to provide the address of the existing buffer -
  the C interface stub will load the registers from this buffer. That is, the
  \emph{snd\_reg} parameter will be: \\ \\
  \hspace*{20pt} {\footnotesize\verb+(l4_ipc_reg_msg_t*) buf +} 
  
\item \emph{If a long message is required, decide on a message format.}
  \\ \\ \hspace*{10pt}
  For this example, the message format will need to include a header,
  a two dword in-line buffer (want to send a 10 dword buffer with first
  8 dwords in registers) and one string dope. Thus, can declare the 
  message as: \\ \\
  \hspace*{20pt} {\footnotesize\verb+typedef struct msg {+} \\
    \hspace*{60pt} {\footnotesize\verb+l4_msghdr_t msghdr;+} \\
    \hspace*{60pt} {\footnotesize\verb+dword_t buf[2];+} \\
    \hspace*{60pt} {\footnotesize\verb+l4_strdope_t strdope;+}\\
    \hspace*{20pt} {\footnotesize\verb+} msg_t;+} \\ \\
  \hspace*{20pt} {\footnotesize\verb+msg_t lmsg;+}

  \hspace*{10pt} The message descriptor is given by: \\ \\
  \hspace*{20pt} {\footnotesize\verb+(msg_t *) (&lmsg & ~(L4_IPC_SHORT_FPAGE | L4_IPC_DECEIT))+} \\ \\
  \hspace*{10pt} The message descriptor is just the address of 
  the message buffer with the last two bits masked out (ie. $m$ =
  0 and $d$ = 0) to indicate a non-deceiving send 
  operation with no mappings.     

  \begin{enumerate}
    
  \item \emph{Fill in the message header.} \\
    
    \hspace*{10pt} The important point to note is that we need
    to send two dwords and one string and and that no fpages are 
    being received. Note also that the \emph{send dope} and the
    \emph{size dope} will be the same since there is no receive 
    involved.\\ \\
    \hspace*{20pt} {\footnotesize\verb+lmsg.msghdr.snd_dope.msgdope = 0; /* zero out all fields */+} \\ 
    \hspace*{20pt} {\footnotesize\verb+lmsg.msghdr.snd_dope.md.dwords = 2; /* number of dwords to send */+} \\
    \hspace*{20pt} {\footnotesize\verb+lmsg.msghdr.snd_dope.md.strings = 1; /* number of strings to send */+}\\
    \hspace*{20pt} {\footnotesize\verb+lmsg.msghdr.size_dope.msgdope = lmsg.msghdr.snd_dope.msgdope; /* send & size dopes */+}\\
    \hspace*{20pt} {\footnotesize\verb+lmsg.msghdr.rcv_fpage.fpage = 0; /* not receiving any fpages */+} \\ \\
    \hspace*{10pt} Note that the the first and last expressions are
    not strictly necessary as the kernel only looks at the
    fields that are actually used.  
    
  \item \emph{Copy in fpage descriptors (if any).} \\ \\
    \hspace*{10pt} No fpages in this example.
    
  \item \emph{Copy in by-value data (if any).} \\ \\
    \hspace*{10pt} No in-memory fpages but have two dwords of 
    by-value data. \\ \\
    \hspace*{20pt} {\footnotesize\verb+lmsg.buf[0] = buf[8];+} \\
    \hspace*{20pt} {\footnotesize\verb+lmsg.buf[1] = buf[9];+} 
    
  \item \emph{Copy in (send) string dopes, one for each string.} \\ \\
    \hspace*{10pt} One string dope describing one send string and 
    no receive string. \\ \\
    \hspace*{20pt} {\footnotesize\verb+lmsg.strdope.snd_size = strlen(sbuf); /* size of string in bytes */+}\\
    \hspace*{20pt} {\footnotesize\verb+lmsg.strdope.snd_str = sbuf; /* start of string */+}\\
    \hspace*{20pt} {\footnotesize\verb+lmsg.strdope.rcv_size = 0; /* no receive string */+}\\
    \hspace*{20pt} {\footnotesize\verb+lmsg.strdope.rcv_str = 0;+} 
  \end{enumerate} 
  
\item \emph{Determine the thread id of the desired receiver thread.} \\
  \\ \hspace*{10pt}
  The thread id of $\sigma_0$ is given by {\footnotesize\verb+SIGMA0_TID+} 
  (constant defined in {\footnotesize\verb+sigma0.h+}). 
  
\item \emph{Determine the desired timeout period.} \\ \\ \hspace*{10pt}
  An infinite timeout period is provided by the constant 
  {\footnotesize\verb+L4_IPC_NEVER+} (defined in {\footnotesize\verb+types.h+}). 
  
\item \emph{Provide the parameters for one of the C interface 
    procedures.} \\ \\ \hspace*{10pt}
  Once again, the most appropriate C interface procedure is:\\ \\
  \hspace*{20pt} {\footnotesize\verb+int l4_mips_ipc_send(l4_thread_t dest,+} \\ 
  \hspace*{130pt} {\footnotesize\verb+const void *snd_msg,+} \\
  \hspace*{130pt} {\footnotesize\verb+l4_ipc_reg_msg_t *snd_reg,+}\\ 
  \hspace*{130pt} {\footnotesize\verb+l4_timeout_t timeout,+}\\ 
  \hspace*{130pt} {\footnotesize\verb+l4_msgdope_t *result)+}\\ \\
  \hspace*{10pt} Filling in the parameters using information from steps
  1-5 gives the desired system call: \\ \\
  \hspace*{20pt} {\footnotesize\verb+l4_mips_ipc_send(SIGMA0_TID,+}\\
  \hspace*{80pt} {\footnotesize\verb+(msg_t *) (((dowrd_t)&lmsg) & ~(L4_IPC_SHORT_FPAGE | L4_IPC_DECEIT)),+}\\
  \hspace*{80pt} {\footnotesize\verb+(l4_ipc_reg_msg_t*) buf, L4_IPC_NEVER, &result);+}
  
\end{enumerate}

%%%%%%%%%%%%%%%%%%%%%%%%%%
%%%%%%%%%%%%%%%%%%%%%%%%%%

\subsubsection{Receiving}

%%%%%%%%%%%%%%%%%%%%%%%%%%
%%%%%%%%%%%%%%%%%%%%%%%%%%


The following is a suggested procedure that can be followed to receive
a message.

\begin{enumerate}

\item Declare a result status variable.

\item Declare a register buffer for the short part of the message.

\item If a long message is required, decide on a message format and
  determine the message descriptor. The format is determined by the
  number of fpages, the amount of by-value data and the number of
  strings to be received.
        
  \begin{enumerate}

  \item Fill in the message header.
  \item Copy in (receive) string dopes, one for each string.

  \end{enumerate}

\item For a closed receive, determine the thread id of the desired
  sender thread. For an open receive, declare a variable to store the
  thread id of the sender.
  
\item Determine the desired timeout period.

\item Provide the parameters for one of the C interface procedures
  which allows an IPC receive.

\end{enumerate}

\example{Receiving a short message}

This example illustrates receiving a short message. Assume the
following conditions:

\begin{itemize}

\item A single non-deceiving IPC receive from any thread.
\item Purely register based message (short message).
\item No timeout (indefinite blocking).

\end{itemize}

Then following the suggested procedure:

\begin{enumerate}

\item \emph{Declare a result status variable.} \\ \\
  \hspace*{20pt} {\footnotesize\verb+l4_msgdope_t result;+}

\item \emph{Declare a register buffer.} \\ \\
  \hspace*{20pt} {\footnotesize\verb+l4_ipc_reg_msg_t rmsg;+} \\ 

\item \emph{If a long message is required, decide on a message format.}
  \\ \\ \hspace*{10pt}
  Not required for a short message. To indicate that the message is 
  purely message based, use {\footnotesize\verb+L4_IPC_SHORT_MSG+} (constant defined in
  {\footnotesize\verb+ipc.h+}) for the message descriptor (\emph{rcv\_msg} parameter).

\item \emph{ For an open receive, declare a variable to store the thread id of the sender.} \\
  \\ \hspace*{20pt} {\footnotesize\verb+l4_threadid_t thrdid;+}
  
\item \emph{Determine the desired timeout period.} \\ \\ \hspace*{10pt}
  An infinite timeout period is provided by the constant 
  {\footnotesize\verb+L4_IPC_NEVER+} (defined in {\footnotesize\verb+types.h+}). 
  
\item \emph{Provide the parameters for one of the C interface 
    procedures.} \\ \\ \hspace*{10pt}
  For the given conditions the most appropriate C interface procedure 
  is:\\ \\
  \hspace*{20pt} {\footnotesize\verb+int l4_mips_ipc_wait(l4_thread_t *src,+} \\ 
  \hspace*{130pt} {\footnotesize\verb+const void *rcv_msg,+} \\
  \hspace*{130pt} {\footnotesize\verb+l4_ipc_reg_msg_t *rcv_reg,+}\\ 
  \hspace*{130pt} {\footnotesize\verb+l4_timeout_t timeout,+}\\ 
  \hspace*{130pt} {\footnotesize\verb+l4_msgdope_t *result)+}\\ \\
  \hspace*{10pt} Filling in the parameters using information from steps
  1-5 gives the desired system call: \\ \\
  \hspace*{20pt} {\footnotesize\verb+l4_mips_ipc_wait(&thrdid, L4_IPC_SHORT_MSG,+}\\
  \hspace*{140pt} {\footnotesize\verb+&rmsg, L4_IPC_NEVER, &result);+}
  
\end{enumerate}  

\example{Receiving a long message}

This example illustrates receiving a long message that contains
in-line as well as string data. Assume the following conditions:

\begin{itemize}

\item A single non-deceiving IPC receive from a specific thread. Assume the sender id has been 
  stored in a variable with the declaration: \\ \\
  \hspace*{10pt} {\footnotesize\verb+ l4_threadid_t senderid;+}
\item Uses in memory message buffer (long message).
\item No fpages.
\item 10 dwords of in-line by-value data is to be received.
\item One string is to be received and stored in a character string buffer with the declaration: \\ \\
  \hspace*{10pt} {\footnotesize\verb+char sbuf[MAX_BUF];+}.\\
  \hspace*{15pt} where {\footnotesize\verb+MAX_BUF+} is a constant.
\item No timeout (indefinite blocking).
  
\end{itemize}

The above assumptions made by the receiver actually defines the
send/receive protocol. In particular, the receiver expects \emph{at
  most} 10 dwords of in-line data and one string of maximum size
{\footnotesize\verb+MAX_BUF+}. Note that the sender may actually send less than the
expected maximum (but not more). That is, there may be less than 10
dwords of in-line data, no string data or a string smaller than the
maximum size.

Then following the suggested procedure:

\begin{enumerate}

\item \emph{Declare a result status variable.} \\ \\
  \hspace*{20pt} {\footnotesize\verb+l4_msgdope_t result;+}
  
\item \emph{Declare a register buffer.} \\ \\
  \hspace*{20pt} {\footnotesize\verb+l4_ipc_reg_msg_t rmsg;+} \\ 
  
\item \emph{If a long message is required, decide on a message format.}
  \\ \\ \hspace*{10pt}
  For this example, the message format will need to include a header,
  a two dword in-line buffer (want to receive a 10 dwords with first
  8 dwords in registers) and one string dope. Thus, can declare the 
  message as: \\ \\
  \hspace*{20pt} {\footnotesize\verb+typedef struct msg {+} \\
    \hspace*{60pt} {\footnotesize\verb+l4_msghdr_t msghdr;+} \\
    \hspace*{60pt} {\footnotesize\verb+dword_t buf[2];+} \\
    \hspace*{60pt} {\footnotesize\verb+l4_strdope_t strdope;+}\\
    \hspace*{20pt} {\footnotesize\verb+} msg_t;+} \\ \\
  \hspace*{20pt} {\footnotesize\verb+msg_t lmsg;+}
  
  \hspace*{10pt} The message descriptor is given by: \\ \\
  \hspace*{20pt} {\footnotesize\verb+(msg_t *) (((dword_t)&lmsg) & ~(L4_IPC_SHORT_FPAGE | L4_IPC_DECEIT))+} \\ \\
  \hspace*{10pt} The message descriptor is just the address of 
  the message buffer with the last two bits masked out (ie. $m$ =
  0 and $d$ = 0) to indicate a non-deceiving receive 
  operation with no mappings.
  
  \begin{enumerate}
    
  \item \emph{Fill in the message header.} \\
    
    \hspace*{10pt} The important point to note is that we need
    to send two dwords and one string and and that no fpages are 
    being received. Note also that the \emph{size dope} will be greater
    than the \emph{send dope} since there is no send 
    involved.\\ \\ 
    \hspace*{20pt} {\footnotesize\verb+lmsg.msghdr.size_dope.md.dwords = 2; /* number of dwords */+} \\
    \hspace*{20pt} {\footnotesize\verb+lmsg.msghdr.size_dope.md.strings = 1; /* number of strings */;+}\\
    \hspace*{20pt} {\footnotesize\verb+lmsg.msghdr.snd_dope = 0; /* no send (actually not needed)*/+}\\
    
  \item \emph{Copy in (receive) string dopes, one for each string.} \\ \\
    \hspace*{10pt} One string dope describing one receive string and 
    no send string. \\ \\
    \hspace*{20pt} {\footnotesize\verb+lmsg.strdope.rcv_size = MAX_BUF; /* max bytes receive string */+}\\
    \hspace*{20pt} {\footnotesize\verb+lmsg.strdope.rcv_str = sbuf; /* start of string */+}\\
    \hspace*{20pt} {\footnotesize\verb+lmsg.strdope.snd_size = 0; /* no send string */+}\\
    \hspace*{20pt} {\footnotesize\verb+lmsg.strdope.snd_str = 0;+} 
    
  \end{enumerate}
  
\item \emph{ For a closed receive, determine the thread id of the desired sender thread. } \\ \\
  \hspace*{10pt} The sender id is stored in the variable {\footnotesize\verb+senderid+}
  
\item \emph{Determine the desired timeout period.} \\ \\ \hspace*{10pt}
  An infinite timeout period is provided by the constant 
  {\footnotesize\verb+L4_IPC_NEVER+} (defined in {\footnotesize\verb+types.h+}). 
  
\item \emph{Provide the parameters for one of the C interface 
    procedures.} \\ \\ \hspace*{10pt}
  For the given conditions the most appropriate C interface procedure 
  is:\\ \\
  \hspace*{20pt} {\footnotesize\verb+int l4_mips_ipc_receive(l4_thread_t src,+} \\ 
  \hspace*{130pt} {\footnotesize\verb+const void *rcv_msg,+} \\
  \hspace*{130pt} {\footnotesize\verb+l4_ipc_reg_msg_t *rcv_reg,+}\\ 
  \hspace*{130pt} {\footnotesize\verb+l4_timeout_t timeout,+}\\ 
  \hspace*{130pt} {\footnotesize\verb+l4_msgdope_t *result)+}\\ \\
  \hspace*{10pt} Filling in the parameters using information from steps
  1-5 gives the desired system call: \\ \\
  \hspace*{20pt} {\footnotesize\verb+l4_mips_ipc_receive(senderid,+}\\
  \hspace*{80pt} {\footnotesize\verb+(msg_t *) ((dword_t)&lmsg) & ~(L4_IPC_SHORT_FPAGE | L4_IPC_DECEIT))+}\\
  \hspace*{80pt} {\footnotesize\verb+&rmsg, L4_IPC_NEVER, &result);+}
  
\end{enumerate}

Note that the kernel does not tell you how much data was actually
transmitted, this information must either be explicitly encoded in the
message or must be implicit in the protocol used between sender and
receiver.

%%%%%%%%%%%%%%%%%%%%%%%%%%
%%%%%%%%%%%%%%%%%%%%%%%%%%

\section{Clans and Chiefs}

%%%%%%%%%%%%%%%%%%%%%%%%%%
%%%%%%%%%%%%%%%%%%%%%%%%%%


%%%%%%%%%%%%%%%%%%%%%%%%%%
%%%%%%%%%%%%%%%%%%%%%%%%%%

\subsection{Concepts}

%%%%%%%%%%%%%%%%%%%%%%%%%%
%%%%%%%%%%%%%%%%%%%%%%%%%%

As described in Section~\ref{s:i:c&c}, the clans and chiefs concept is
one of the \emph{security mechanisms} used by L4 to allow
\emph{protection policies} to be implemented for its IPC message
transfers.

Clans are created via a hierarchy of tasks. A task that creates
another task becomes the chief of that task and the set of tasks
created by the chief forms the chief's clan. Note that there is an
implementation limit on the depth of the task hierarchy\footnote{16 on
  R4k.}.

Security is provided through a number of restrictions and enforced
protocols:

\begin{itemize}

\item A task can only (directly) kill another task if that task is in
  its own clan or (indirectly) by killing its chief.

\item Intra-clan messages (those that don't cross clan boundaries) are
  delivered directly from sender to receiver.

\item Inter-clan messages (those that cross clan boundaries), whether
  incoming or outgoing, are routed from the sender to the sender's
  chief instead of going to the specified receiver (in the other
  clan). The chief has access to all parts of the message and can
  inspect and modify the message before forwarding it or even suppress
  the message entirely. If forwarding the message, the chief can make
  use of deceiving IPC so that the message will seemingly come from
  the original sender.

\item Deceiving is \emph{direction-preserving}: L4 will only allow the
  deceit if, on an outgoing message, the virtual sender ID belongs
  (directly or indirectly) to the sender's clan, or, on a message
  going inside the clan the virtual sender is external to the sender's
  clan. Furthermore, the receiver is alerted to the deceit (via the
  \(d\)-bit in the condition code returned from the call, see
  Section~\ref{s:ccode}).


\end{itemize}


%%%%%%%%%%%%%%%%%%%%%%%%%%
%%%%%%%%%%%%%%%%%%%%%%%%%%

\subsection{Usage and Cost}

%%%%%%%%%%%%%%%%%%%%%%%%%%
%%%%%%%%%%%%%%%%%%%%%%%%%%

Some usage examples of clans and chiefs are:

\begin{description}

\item[network proxies] Have a master clan for each node which receives
  all inter-node IPC and forwards it.

\item[system upgrade] Have a clan for each \emph{system version}. The
  chief for each clan can then translate message formats between old
  and new versions as well as any other necessary adjustments.

\item[user-based access permissions] All tasks for a user are
  encapsulated in a single clan. The chief of each clan (user) can
  implement arbitrary access protocols.

\item[confinement] Encapsulating a task in its own clan may be used to
  prevent it from leaking data or mounting denial-of-service attacks.

\end{description}

The cost of \emph{supporting} clans and chiefs is about four cycles
per IPC.  The cost of \emph{using} clans and chiefs is the product of
the number of IPC operations and the number of chiefs involved. This
last point implies that the task hierarchy should be kept as flat as
possible in the interests of efficiency. \\ 

\example{Clans and chiefs}

This example illustrates one possible inter-clan IPC scenario. It also
demonstrates the use of deceiving IPC. Figure~\ref{fig:clanex} shows
the setup for this example. Task $T_2$ is the chief of task $T_1$
while task $T_3$ is the chief of task $T_4$. $T_1$ wants to send an
inter-clan message to $T_4$. The dashed line shows the virtual path of
the message from $T_1$ to $T_4$ but the actual path is shown by the
normal line. $T_1$ sends to $T_4$ but the message is intercepted by
its chief, $T_2$. $T_2$ tries to pass the message to $T_4$ but the
message is actually intercepted by $T_3$ which finally forwards the
message to $T_4$. Deceiving IPC is used when $T_2$ and $T_3$ forwards
their messages. The effect of this is that $T_4$ will see the message
it eventually receives as coming directly from $T_1$. \\ 

\begin{figure}[htb]
\begin{center}
\leavevmode
\epsfxsize = .7\textwidth
\epsfbox{eps/clanex.eps}
\end{center}
\caption{Clans and Chiefs Example}
\label{fig:clanex}
\end{figure}

There are a few things to note before proceeding with this example.

\begin{itemize}

\item The use of deceiving IPC as in this example is \emph{not} a
  mandatory part of using the clans and chiefs mechanism. Depending on
  the situation, the user may decide that $T_4$ does not need to know
  that the message was originally from $T_1$ or the identity of the
  original sender may be encoded as part of the message. In either
  case, deceiving IPC need not be used.

  However, deceiving is necessary to maintain proper RPC semantics in
  the case of redirection: If \(T_1\) attempts an RPC-like
  communication with \(T_4\) (using {\footnotesize\verb+l4_mips_ipc_call+}), the
  reply must come from \(T_4\) for \(T_1\) being able to receive it.
  See also Section~\ref{s:u:osstruct} for an example.

\item In terms of sending an inter-clan message and the actual path
  shown in figure~\ref{fig:clanex}, the only communication which must
  take place is from $T_1$ to $T_2$ (i.e.\ from the original sender to
  its chief). The path from $T_2$ to $T_3$ and from $T_3$ to $T_4$ may
  never occur because these chiefs may decide to suppress the message.

\item The thread \(T_2\) actually receiving the intercepted IPC is
  lthread~0 of the chief task.

\end{itemize}

Code fragments for each of the tasks may look as follows.

\begin{description}

\item[$T_1$ code:]  
  $T_1$ does not need to do anything special. It just attempts to
  send directly to $T_4$. So the code is: \\ \\
  \hspace*{10pt} {\footnotesize\verb+l4_ipc_reg_msg_t msg;+} \\
  \hspace*{10pt} {\footnotesize\verb+l4_msgdope_t result;+} \\
  \hspace*{10pt} {\footnotesize\verb+l4_threadid_t t4id;+} \\ \\
  \hspace*{10pt} {\footnotesize\verb+t4id = +} \dots {\footnotesize\verb+ /* assign t2's thread id */+}\\
  \hspace*{10pt} {\footnotesize\verb+l4_mips_ipc_send(t4id, L4_IPC_SHORT_MSG, &msg, L4_IPC_NEVER, &result);+} \\
  \hspace*{40pt} \vdots   
  
\item[$T_2$ code:]
  $T_2$ waits to receive (actually intercept) a message from a client 
  ($T_1$ in this case)
  and passes it on to $T_4$ (after possibly inspecting and modifying 
  the message) using a deceiving send with $T_1$ as the virtual sender. 
  So the code is: \\ \\
  \hspace*{10pt} {\footnotesize\verb+l4_ipc_reg_msg_t msg;+} \\
  \hspace*{10pt} {\footnotesize\verb+l4_msgdope_t result;+} \\
  \hspace*{10pt} {\footnotesize\verb+l4_threadid_t clientid, t4id;+} \\ \\
  \hspace*{10pt} {\footnotesize\verb+l4_mips_ipc_wait(&clientid, L4_IPC_SHORT_MSG, &msg, L4_IPC_NEVER, &result);+} \\ \\
  \hspace*{10pt} {\footnotesize\verb+/* Inspect/modify message before forwarding */+}\\ \\
  \hspace*{10pt} {\footnotesize\verb+t4id = +} \dots {\footnotesize\verb+ /* assign t4's thread id */+}\\ \\
  \hspace*{15pt} {\footnotesize\verb+/* Send deceiving IPC to T4 with client as virtual sender */+} \\
  \hspace*{10pt} {\footnotesize\verb+l4_mips_ipc_send_deceiving(t4id, clientid, L4_IPC_SHORT_MSG, &msg,+} \\
  \hspace*{10pt} {\footnotesize\verb+                           L4_IPC_NEVER, &result);+} \\ 
  \hspace*{40pt} \vdots   
  
\item[$T_3$ code:]
  $T_3$ intercepts the message from $T_2$ and directly forwards its 
  message to $T_4$. Note that $T_3$ believes it receives its message 
  directly from $T_1$ because of the deceiving IPC used by $T_2$.
  The code for $T_3$ is similar to that for $T_2$ and is not repeated.
  
\item[$T_4$ code:]
  $T_4$ does not need to know of the actual path taken by the message.
  It simply receives the message and thinks (ignoring the
  \(d\)-bit in the returned condition code) that it has come directly
  from $T_1$. So the code is: \\ \\
  \hspace*{10pt} {\footnotesize\verb+l4_ipc_reg_msg_t msg;+} \\
  \hspace*{10pt} {\footnotesize\verb+l4_msgdope_t result;+} \\
  \hspace*{10pt} {\footnotesize\verb+l4_threadid_t senderid;+} \\ \\
  \hspace*{10pt} {\footnotesize\verb+l4_mips_ipc_wait(&senderid, L4_IPC_SHORT_MSG, &msg, L4_IPC_NEVER, &result);+} \\
  \hspace*{40pt} \vdots   
  
\end{description}

This code assumes that $T_2$ and $T_3$ know the destination of the IPC
(it could be explicitly encoded in the message or implicitly in the
protocol). Practically this approach is most useful in such cases as a
server's using a dedicated receiver thread which forwards requests to a
number of worker threads which then reply directly back to the client
(using deceiving IPC).

A more appropriate way for a chief dealing with redirected IPC goes as
follows:

\begin{description}
\item[$T_2$ code:] ~\\
  \hspace*{10pt} {\footnotesize\verb+l4_ipc_reg_msg_t msg;+} \\
  \hspace*{10pt} {\footnotesize\verb+l4_msgdope_t result;+} \\
  \hspace*{10pt} {\footnotesize\verb+l4_threadid_t clientid, destid;+} \\ \\
  \hspace*{10pt} {\footnotesize\verb+l4_mips_ipc_chief_wait(&clientid, &destid, L4_IPC_SHORT_MSG, &msg,+} \\
  \hspace*{10pt} {\footnotesize\verb+                       L4_IPC_NEVER, &result);+} \\ \\
  \hspace*{10pt} {\footnotesize\verb+/* Inspect/modify message before forwarding */+}\\ \\
  \hspace*{15pt} {\footnotesize\verb+/* Send deceiving IPC to intended destination with client as virtual sender */+} \\
  \hspace*{10pt} {\footnotesize\verb+l4_mips_ipc_send_deceiving(destid, clientid, L4_IPC_SHORT_MSG, &msg,+} \\
  \hspace*{10pt} {\footnotesize\verb+                           L4_IPC_NEVER, &result);+} \\ 
  \hspace*{40pt} \vdots   
  
\end{description}



% 
%%% Local Variables: %%%
%%% mode: latex %%%
%%% End: %%%
		% updated 11/8/99

% Permission to make digital/hard copy of part or all of this work for
% personal or classroom use is granted without fee provided that copies
% are not made or distributed for profit or commercial advantage, the
% copyright notice, the title of the publication and its date appear, and
% notice is given that copying is by permission of the authors. To copy
% otherwise, to republish, to post on servers, or to redistribute to lists
% requires prior specific permission and/or a fee.
%
% Copyright (C) 1998 by Gernot Heiser, The University of New South Wales.


% L4 user manual non-IPC system calls chapter
% Author: Alan Au
% Created On: 5 January 1998
% Last Modified: 29 May 1998

%%%%%%%%%%%%%%%%%%%%%%%%%%
%%%%%%%%%%%%%%%%%%%%%%%%%%

\chapter{\label{c:other}Other L4 System Calls}

%%%%%%%%%%%%%%%%%%%%%%%%%%
%%%%%%%%%%%%%%%%%%%%%%%%%%


%%%%%%%%%%%%%%%%%%%%%%%%%%
%%%%%%%%%%%%%%%%%%%%%%%%%%

\section{Task Creation and Deletion}

%%%%%%%%%%%%%%%%%%%%%%%%%%
%%%%%%%%%%%%%%%%%%%%%%%%%%

\label{sec:tasks}

A task can be in either an \emph{active} or an \emph{inactive} state.
Creating an active task creates a new address space as well as the
maximum number of threads for a task (128). Initially, all the threads
of an active task except for one (called lthread 0) are
\emph{inactive}. In contrast, an inactive task has no address space
and no threads (whether active or not) and thus consumes no resources.

The kernel only allows a certain number of tasks to be created on a
first-come-first-served basis with subsequent attempts to create a new
task failing (see section~\ref{sec:ralloc}). Thus, the purpose of
creating an inactive task is essentially to reserve the right to create
an active task. Further, an inactive task can be donated to another
chief which effectively transfers the right to create an active
task. The chief of a new active task is the task that created it. A task
created as inactive can have a new chief (\emph{task donation}).

Deleting a task removes the address space and all associated threads of
the task. A task can only be deleted by its chief or indirectly by (a
chief higher up in the task hierarchy) deleting its chief.

Task creation and deletion is done by using the {\footnotesize\verb+task_new+} system
call. This system call first deletes a (active or inactive) task and
creates a new (active or inactive) one. If the task is created active
it gets the same task number (as provided in the \emph{dest}
parameter) but a different version number, hence producing a different
ID. An active task is created by providing a valid pager id to the
system call while a nil pager produces an inactive task.

Note that there is no separate task deletion system call as such. To
kill a task, simply create a new inactive task providing the id of the
task to be killed to the {\footnotesize\verb+task_new+} system call (as the
\emph{dest} parameter).

Creating an \emph{active} task requires the caller to supply
\begin{itemize}
\item a pager (the pager for the new task's lthread 0, as well as the
default pager for all further threads),
\item a start address and an initial stack pointer
\item a maximum scheduling priority (MCP), and
\item an exception handler (the handler for the new task's lthread 0, as
well as the default handler for all further threads).
\end{itemize}
Creating an \emph{inactive} task requires the caller to specify
\begin{itemize}
\item a null pager (which identifies the new task as being inactive),
and
\item a new chief for the task (which can be the same as the calling
task, if not then the call donates the task to the new chief).
\end{itemize}
An example of task creation can be found in Section~\ref{s:ex-task}.

%%%%%%%%%%%%%%%%%%%%%%%%%%
%%%%%%%%%%%%%%%%%%%%%%%%%%

\section{Thread Related System Calls}

%%%%%%%%%%%%%%%%%%%%%%%%%%
%%%%%%%%%%%%%%%%%%%%%%%%%%

%%%%%%%%%%%%%%%%%%%%%%%%%%
%%%%%%%%%%%%%%%%%%%%%%%%%%

\subsection{Thread Manipulation}
\label{sec:exregs}

%%%%%%%%%%%%%%%%%%%%%%%%%%
%%%%%%%%%%%%%%%%%%%%%%%%%%

As mentioned already, an active task is created with a full set (128) of
threads but with only one thread active (lthread 0). A thread is
activated by setting its instruction pointer (IP) and stack pointer (SP)
to valid values. Once active, a thread cannot be deactivated (other than
by deleting its task). To stop a thread from running it needs to be
blocked on an IPC which will never succeed.

A thread's instruction and stack pointers, along with its exception
handler and pager, are set by manipulating the thread's register
values through the {\footnotesize\verb+lthread_ex_regs+} system call. Providing the
invalid value (-1) for any of these to this system call will retain the old
values. {\footnotesize\verb+lthread_ex_regs+} also gives back the old values of the
instruction pointer, stack pointer, exception handler and pager. Thus,
the call can also be used to perform a (logical) thread switch
(exchange registers of running thread with saved a one, this supports
thread management by an OS personality ) as well as save a thread's
current context (by providing invalid values only). A thread's pending
IPC's are cancelled and those in progress are aborted by this system
call.

%%%%%%%%%%%%%%%%%%%%%%%%%%
%%%%%%%%%%%%%%%%%%%%%%%%%%

\subsection{Release CPU}

%%%%%%%%%%%%%%%%%%%%%%%%%%
%%%%%%%%%%%%%%%%%%%%%%%%%%

A thread can use the {\footnotesize\verb+thread_switch+} system call to voluntarily
release the CPU. The releasing thread can specify a specific thread to
which to donate its remaining time slice. If ready, the thread
receiving the donation obtains the remaining time of the other thread
on top of its own time slice. Alternatively, if the receiving thread
is not ready or if the releasing thread does not specify a destination
thread as part of the system call, the caller's remaining time slice
is simply forfeited and normal scheduling takes place immediately.

An important use of this system call is to implement user-level
scheduling. The user-level scheduler is (the only) thread running at
highest priority, so it will be run by the \micro-kernel whenever the
kernel scheduler is invoked. The user-level scheduler then selects the
next thread it wants to run and donates its time slice to it.

%%%%%%%%%%%%%%%%%%%%%%%%%%
%%%%%%%%%%%%%%%%%%%%%%%%%%

\subsection{Thread Scheduling}

%%%%%%%%%%%%%%%%%%%%%%%%%%
%%%%%%%%%%%%%%%%%%%%%%%%%%

Thread scheduling in L4 is controlled by three parameters:
\emph{timeslice length, thread priority and maximum controlled
  priority (MCP)}.

\subsubsection{Timeslice Length}

Each L4 thread has a timeslice length associated with it.  The timeslice value
can range from 0  to {\footnotesize\verb+MAX_TIMESLICE+} and can  vary between
individual threads  of  a task. Each  thread  is scheduled for  the  timeslice
length currently associated with it. When a thread's time quantum expires, the
scheduler  selects the next runnable  thread  as  described  in the  following
section.

A timeslice length of zero is valid.  A thread with zero timeslice is
taken out of the ready queue and therefore never scheduled (until it
is given a non-zero timeslice length).

Note that a thread's timeslice length is in no way determined by its
priority. It is valid for threads of the same priority to have
different timeslice lengths. A thread initially gets the same
timeslice length as its parent and that value can only be changed via
a {\footnotesize\verb+l4_thread_schedule+} system call.

\subsubsection{Priority}

The kernel defines 256 levels of priority in the range
[255..0] with 255 being the highest priority and 0 the lowest. L4's
internal scheduler uses multiple-level round-robin queues such that
there is a queue (possibly empty) associated with each priority level.
All the queues taken together form the kernel's ready queue.

Figure~\ref{fig:pbl} shows an example of a ready queue in L4. Each circle 
represents a thread in a particular round robin queue. 

\begin{figure}[htb]
\begin{center}
\leavevmode
\epsfxsize = .7\textwidth
\epsfbox{eps/priority.eps}
\end{center}
\caption{Example Ready Queue}
\label{fig:pbl}
\end{figure}

Each thread has an associated priority at any given time. The priority
determines which round robin queue the thread belongs to in the kernel
ready queue.  Changing a thread's priority (via
{\footnotesize\verb+l4_thread_schedule+}) will change the queue it belongs to.

L4 priorities are absolute.  On each scheduling event, the scheduler
will always select the next thread to run from the head of the highest
priority queue that is currently non-empty. For example, in
Figure~\ref{fig:pbl}, the scheduler would take the thread at the head
of the queue associated with priority 254.

\subsubsection{Maximum Controlled Priority}
\label{sec:mcp}

Unlike the timeslice length and timeslice priority, the MCP is not
thread based but rather task based. The MCP of a task is specified at
creation time (see Section~\ref{sec:tasks}) and all threads in the
task will share this MCP value.

Any thread can change another thread's scheduling parameters (timeslice
length and priority) by invoking L4's {\footnotesize\verb+l4_thread_schedule+} system
call under the right conditions. {\footnotesize\verb+l4_thread_schedule+} works if
and only if src.mcp $\geq$ dest.prio AND src.mcp $\geq$ new prio.
Otherwise, it will not change the destination's status but its effects
on src are undefined.

\subsubsection{Scheduling Parameter Inheritance}

When an L4 task is created (by calling {\footnotesize\verb+l4_task_new+}), only the
MCP is specified but neither the time slice length nor priority is
given. Simlarily, creating an L4 thread (by calling
{\footnotesize\verb+l4_thread_ex_regs+}) does not explicitly require any of the
scheduling parameters to be provided. Thus there are implicit
scheduling parameter inheritance rules defined for new tasks and
threads.

In the following description of the inheritance rules $new~x$
represents the value of $x$ given as a parameter in the relevant
system call ({\footnotesize\verb+l4_task_new+} or {\footnotesize\verb+l4_thread_ex_regs+}),
$src.x$ represents the value of $x$ in the creator task/thread and
$dest.x$ represents the value of $x$ in the task/thread being
created. $x$ is one of either $mcp$, $tsl$ or $prio$. 

\begin{description}

\item[New task] Only lthread~0 of the new task has its scheduling
  parameters defined.  All other threads will have their scheduling
  parameters defined when they are first activated (except for MCP
  which is task based). So the actual inheritance rules for lthread~0
  of the new task are:

  \begin{itemize}

    \item $dest.mcp = \min\{new~mcp, src.mcp\}$

    \item $dest.tsl = src.tsl$

    \item $dest.prio = src.prio$
      
  \end{itemize}

\item[New thread] In this context, a thread is only considered new
  when it is first ``created''. Recall from Section~\ref{sec:exregs}
  that the system call {\footnotesize\verb+l4_thread_ex_regs+} can be used to
  activate an inactive thread, to do a logical thread switch or just
  to get a thread's state. For the current purposes, only the first
  use of {\footnotesize\verb+l4_thread_ex_regs+} is considered as ``creating'' a new
  thread. This distinction is important because the following
  inheritance rules only apply when a new thread is created.
  \emph{There is no change in scheduling parameters in the other two
    cases.}
  
  \begin{itemize}

  \item $dest.mcp = src.mcp$\footnote{MCP is actually only defined for
      lthread~0 with all other threads in the task using the value from
      lthread~0.}

  \item $dest.tsl = src.tsl$

  \item $dest.prio = src.prio$
      
  \end{itemize}

\end{description}

One final issue must be considered as part of scheduling parameter
inheritance. At boot time, $\sigma_0$ (Section~\ref{sec:sigma0})
starts up any programs that have been marked as initial servers in the
kernel boot image (see Section~\ref{sec:bootstrap} for information on
the L4 bootstrap and Section~\ref{sec:dit} for the boot image).
$\sigma_0$ has maximum values for all its scheduling parameters and
also calls {\footnotesize\verb+l4_task_new+} (to create the initial servers) with the
MCP parameter set to maximum.  Hence all initial servers will have
maximum values for their MCP and thread priority. This behaviour is
sensible because initial servers should form the OS and thus should be
given maximum privileges.


\subsubsection{L4 Scheduling System Call}

The {\footnotesize\verb+l4_thread_schedule+} system call allows setting (kernel)
scheduling parameters of user threads.  It also returns thread states,
in particular accumulated CPU time. The following restrictions and
points must be remembered when calling {\footnotesize\verb+l4_thread_schedule+}:

\begin{itemize}
  
\item For the call to be effective, the MCP condition described above
  must be satisfied.

\item The call cannot increase the destination thread's priority over
  the caller task's own MCP. That is, the value of \emph{prio} given
  in the \emph{param} parameter must not exceed the caller's MCP.

\item The new timeslice length of the destination can be set to any
  value\footnote{Obviously, the timeslice length can only be specified
  with the accuracy allowed by the granularity of the timeslice
  format. Moreover, not all representable values may be possible in the
  particular kernel, and the set of possible values may depend on the
  priority of the target thread. L4 will round the specified value to
  the nearest possible value.} within the interval [0,
  {\footnotesize\verb+MAX_TIMESLICE+}]. One way to set a thread to have
  {\footnotesize\verb+MAX_TIMESLICE+} is to give the maximum value that
  can be specified in timeslice format as the new timeslice length.

\end{itemize}

%%%%%%%%%%%%%%%%%%%%%%%%%%
%%%%%%%%%%%%%%%%%%%%%%%%%%

\subsection{Obtaining Thread Identifiers}

%%%%%%%%%%%%%%%%%%%%%%%%%%
%%%%%%%%%%%%%%%%%%%%%%%%%%

The {\footnotesize\verb+id_nearest+} system call returns the ID of the thread that
would \emph{really} receive a message sent to a specific thread. In
particular, if the destination is:

\begin{description}

        \item[outside invoker's clan:] Return ID of invoker's own chief.

        \item[in same clan as invoker, or inside invoker's clan:] Return ID of destination thread.

        \item[in subclan of invoker's clan:] Return ID of chief (within own
        clan) of topmost subclan.

        \item[nil:] Return invoker's own thread ID. (This is the closest
        thing to a ``null system call'' in L4 and is frequently used for
        benchmarking.)

\end{description}

\begin{figure}[h]
\centering
\mbox{\subfigure[Outside]
        {\epsfig{figure=eps/id1.eps,width=.32\textwidth}}\quad
      \subfigure[Inside]
        {\epsfig{figure=eps/id2.eps,width=.25\textwidth}}\quad
      \subfigure[In Subclan]
        {\epsfig{figure=eps/id3.eps,width=.25\textwidth}}}
\caption{id\_nearest}
\label{fig:id}
\end{figure}

Figure~\ref{fig:id} shows the first three cases. Clan boundaries are
shown as ovals and tasks are shown as circles. An arrow is drawn from
the sender to the destination and the shaded circle is the thread ID
that is returned by {\footnotesize\verb+id_nearest+} when invoked by the sender.

%%%%%%%%%%%%%%%%%%%%%%%%%%
%%%%%%%%%%%%%%%%%%%%%%%%%%

\section{Revoking Mappings}

%%%%%%%%%%%%%%%%%%%%%%%%%%
%%%%%%%%%%%%%%%%%%%%%%%%%%

Mappings can be recalled by using the {\footnotesize\verb+fpage_unmap+} system
call.  The invoker  specifies an fpage to  be revoked from  all address spaces
into which the  invoker mapped directly or indirectly.   The unmapping can  be
done partially  (revert to read-only) or completely  (pages no  longer part of
the other   address  spaces).   As part  of  the  unmapping,  the  invoker can
optionally elect to remove the pages from its own address space.

\cbstart

The {\footnotesize\verb+fpage_unmap+} takes two paramters:

\begin{itemize}

\item \emph{fpage:} This the fpage to  be unmapped subject  to the map mask in
the second  parameter. As with mapping and   granting, the fpage specification
designates all valid mapping  within the address  space region defined  by the
fpage.

\item \emph{map\_mask:} This determines how the unmap is performed
by indicating:

	\begin{enumerate}

	\item The unmap operation - set fpage to read-only
	({\footnotesize\verb+L4_FP_REMAP_PAGE+}) or completely unmap fpage
	({\footnotesize\verb+L4_FP_FLUSH_PAGE+}).

	\item The unmap extent - apply the unmap operation to all other address
	spaces in which the fpage has been mapped but not the original flex
	page ({\footnotesize\verb+L4_FP_OTHER_SPACES+}) or apply the unmap
	operation in every address space including the original
	({\footnotesize\verb+L4_FP_ALL_SPACES+})

	\end{enumerate}

	Note that the  unmap operation and  unmap extent are orthogonal and so
	both should  be  specified (by combining the    two attributes with  a
	logical   OR). The  table   below shows   all  the  valid  values  for
	\emph{map\_mask}.

	\begin{tabular}{|c|c|}  \hline \hline

	\emph{map\_mask} & Description \\ \hline
	
	{\footnotesize\verb+L4_FP_REMAP_PAGE |+} & 
	Map fpage read-only  in all other address spaces\\ 

	{\footnotesize\verb+L4_FP_OTHER_SPACES+} & 
	in which the fpage has been mapped\\ \hline

	{\footnotesize\verb+L4_FP_FLUSH_PAGE |+} & 
	Completely unmap the fpage in all other address spaces\\

	{\footnotesize\verb+L4_FP_OTHER_SPACES+} & 
	in which the fpage has been mapped \\ \hline

	{\footnotesize\verb+L4_FP_REMAP_PAGE |+} & 
	Map fpage read-only \\

	{\footnotesize\verb+L4_FP_ALL_SPACES+} & 
	in all address spaces\\ \hline

	{\footnotesize\verb+L4_FP_FLUSH_PAGE |+} & 
	Completely unmap the fpage \\

	{\footnotesize\verb+L4_FP_ALL_SPACES+} & 
	in all address spaces\\ \hline

	\end{tabular}

\end{itemize}

\cbend


%%%%%%%%%%%%%%%%%%%%%%%%%%
%%%%%%%%%%%%%%%%%%%%%%%%%%

\section{An Example}

%%%%%%%%%%%%%%%%%%%%%%%%%%
%%%%%%%%%%%%%%%%%%%%%%%%%%


\example{Thread and task creation}

This example illustrates how a thread can be made active and how to create an active (sub)task. First a pager thread will be started and then a sub-task will be created using the newly started pager.

Assume we have the following declarations: \\ \\
\hspace*{10pt} {\footnotesize\verb+#define PAGERSTACKSIZE 1024+} \\
\hspace*{10pt} {\footnotesize\verb+dword_t pager_stack[PAGERSTACKSIZE];+} \\
\hspace*{10pt} {\footnotesize\verb+extern void pager(void); /* pager function */+} \\
\hspace*{10pt} {\footnotesize\verb+extern void subtask(void); /* subtask function */+} \\
\hspace*{10pt} {\footnotesize\verb+dword_t oip, osp; /* old insturction & stack pointers */+} \\
\hspace*{10pt} {\footnotesize\verb+l4_thread_id_t excpt, page; /* exception handler and pager */+} \\
\hspace*{10pt} {\footnotesize\verb+l4_thread_id_t pagerid, subtaskid, myid;+} \\

Now, to start a new thread the following steps are needed.

\begin{enumerate}

\item Get the id of own task. \\ \\
  \hspace*{10pt} {\footnotesize\verb+myid = l4_myself();+} \\ \\
  Note that {\footnotesize\verb+l4_myself+} is just a C interface procedure which
  uses the {\footnotesize\verb+id_nearest+} L4 system call with a \emph{nil} parameter.
  
\item Obtain an id for the new thread. Here we just use an
  lthread number one higher than that of the calling thread. (This
  one better be unused!) \\ \\
  \hspace*{10pt} {\footnotesize\verb+ pagerid = myid; /* same task */+} \\
  \hspace*{10pt} {\footnotesize\verb" pagerid.id.lthread += 1; /* new thread */"}

\item For this example, there is to be no change in exception handler
  and pager. That is, the new thread is to have the same exception
  handler and pager as the thread that started it. Recall that this
  is done by providing an invalid (-1) value for the exception handler
  and pager parameters in the {\footnotesize\verb+l4_thread_ex_regs+} system call. \\ \\
  \hspace*{10pt} {\footnotesize\verb+excpt.ID = -1LL;+} \\
  \hspace*{10pt} {\footnotesize\verb+page.ID = -1LL;+}
  
\item Determine the new instruction pointer for the thread. In this
  case, it is just the entry address of our pager function (i.e.
  {\footnotesize\verb+pager+}).

\item Determine the new stack pointer for the thread. We can use any
  region of memory (in this task's address space) reserved for such a
  purpose. The stack starts from high memory and ``grows'' downwards.
  Therefore out stack pointer is\\ \\
  \hspace*{10pt} {\footnotesize\verb+&pager_stack[PAGERSTACK - 1]+}
  
\item Use {\footnotesize\verb+l4_thread_ex_regs+} to start the new thread. \\ \\
  \hspace*{10pt} {\footnotesize\verb+l4_thread_ex_regs(pagerid,+} \\
  \hspace*{105pt} {\footnotesize\verb+(dword_t) pager,+}\\
  \hspace*{105pt} {\footnotesize\verb+(dword_t) &pager_stack[PAGERSTACK - 1],+}\\
  \hspace*{105pt} {\footnotesize\verb+&excpt,+}\\
  \hspace*{105pt} {\footnotesize\verb+&page,+}\\
  \hspace*{105pt} {\footnotesize\verb+&oip,+}\\
  \hspace*{105pt} {\footnotesize\verb+&osp);+}
\end{enumerate}

The second part of this example is to create a new task with the newly
started thread as the pager for the new task. The steps may be as
follows:\label{s:ex-task}

\begin{enumerate}

\item Obtain a new task id. This can be done simply by picking
  any unused task id, here we assume that the caller's task is the
  last one in use so far. In general we need to manage
  task IDs. \\ \\
  \hspace*{10pt} {\footnotesize\verb+subtaskid.ID = myid.ID;+} \\
  \hspace*{10pt} {\footnotesize\verb"subtaskid.id.task += 1; /* just pick an id higher than self */"} \\
  
\item Determine the exception handler for the new task. For this
  example, the exception handler will be set to be the thread that
  creates the new task (i.e. {\footnotesize\verb+myid+}).
  
\item Determine the pager for the new task. For this example, the
  pager will be the previously started thread (i.e. {\footnotesize\verb+pagerid+}).
  
\item Determine the initial instruction pointer for the new task.  In
  this case, it is just {\footnotesize\verb+subtask+}.
  
\item Determine the maximum scheduling priority of the new task. Here
  we assume that this is already set in a variable {\footnotesize\verb+mcp+}.
  
\item Determine the initial stack pointer for the new task. Note the
  difference between this stack pointer and that for creating a new
  thread. A new thread has the same address space as the thread that
  started it (because they belong to the same task) but the new task
  will have a totally new address space. By convention, the stack for
  a new task is placed in the high end of memory though the exact
  place is up to the user. For this example, we place the stack at the
  top end of user addressable memory (address \(2^{40}\) on the R4x00).

\item Use {\footnotesize\verb+l4_task_new+} to create the new task. \\ \\ 
  \hspace*{10pt} {\footnotesize\verb+subtaskid = l4_task_new(subtaskid, mcp, 1<<40, &subtask, pagerid, myid);+}

\end{enumerate}



% 
%%% Local Variables: %%%
%%% mode: latex %%%
%%% End: %%%
	% updated 11/8/99

% Permission to make digital/hard copy of part or all of this work for
% personal or classroom use is granted without fee provided that copies
% are not made or distributed for profit or commercial advantage, the
% copyright notice, the title of the publication and its date appear, and
% notice is given that copying is by permission of the authors. To copy
% otherwise, to republish, to post on servers, or to redistribute to lists
% requires prior specific permission and/or a fee.
%
% Copyright (C) 1998 by Gernot Heiser, The University of New South Wales.


% L4 user manual: How L4 is used
% Author: Alan Au
% Created On: 20 December 1997
% Last Modified: 1 June 1998

%%%%%%%%%%%%%%%%%%%%%%%%%%
%%%%%%%%%%%%%%%%%%%%%%%%%%

\chapter{\label{c:use}Using L4}

%%%%%%%%%%%%%%%%%%%%%%%%%%
%%%%%%%%%%%%%%%%%%%%%%%%%%

%%%%%%%%%%%%%%%%%%%%%%%%%%
%%%%%%%%%%%%%%%%%%%%%%%%%%

\section{Bootstrap}
\label{sec:bootstrap}

%%%%%%%%%%%%%%%%%%%%%%%%%%
%%%%%%%%%%%%%%%%%%%%%%%%%%

At boot time, the kernel image is loaded into resident memory (RAM).
The first process started (in kernel mode) by the L4 bootstrap is
$\sigma_0$ (see Section~\ref{sec:sigma0}). $\sigma_0$ then in turn
starts up all the \emph{initial servers} in user mode. All initial
servers have $\sigma_0$ registered as their pager and exception
handler (i.e.$\sigma_0$ is their chief). $\sigma0$ does not allocate
any stack space for the initial servers. The initial servers must
allocate their own stack space and point the variable \emph{\_sp} to
the beginning of that allocated space.

Kernel images are built by a bootstrap linker called DIT (see
Section~\ref{sec:dit}). In building the image, the DIT identifies the
initial servers by marking them as such in the boot image. The boot
image layout is then:

\begin{figure}[h]
\begin{center}
\leavevmode
\epsfxsize = .8\textwidth
\epsfbox{eps/l4-dit.eps}
\end{center}
\caption{L4/MIPS Boot Image}
\label{fig:boot}
\end{figure}

%%%%%%%%%%%%%%%%%%%%%%%%%%
%%%%%%%%%%%%%%%%%%%%%%%%%%

\section{Resource Allocation}

%%%%%%%%%%%%%%%%%%%%%%%%%%
%%%%%%%%%%%%%%%%%%%%%%%%%%

\label{sec:ralloc}

L4 manages several resources which it allocates on a
first-come-first-served basis. This means that the first task to
request for a particular resource will be granted it while all
subsequent requests for the same resource will fail. These resources
can only be requested by initial servers (i.e.\ by tasks which are
direct members of \(\sigma_0\)'s clan).

Specifically, the resources with this allocation policy are:

\begin{description}

\item[physical frames] All physical frames are mapped idempotently
  (i.e.\ virtual address = physical address) by $\sigma_0$. Except for
  some special kernel reserved memory, $\sigma_0$ only gives each
  frame once to the first task asking for it (see
  Section~\ref{sec:sigma0}).

\item[device addresses] Device access is memory mapped. Like physical
  memory, Pages containing device addresses are mapped once to the
  first requester.

\item[interrupts] Each interrupt can have at most one interrupt
  handler. A thread installs itself as this handler by following a set
  protocol (see Section~\ref{sec:interrupts}). The first thread to do
  so for a particular interrupt becomes the handler while all
  subsequent attempts by other threads will fail.

\item[tasks] Upon system initialisation, the full set of tasks (2048
  on R4k) is created but in an \emph{inactive} state. Tasks are
  acquired using L4's {\footnotesize\verb+task_new+} system call. Once all available
  tasks have been given out, no more tasks can be asked for (see
  Section~\ref{sec:tasks}).

\end{description}

%%%%%%%%%%%%%%%%%%%%%%%%%%
%%%%%%%%%%%%%%%%%%%%%%%%%%

\section{$\sigma_0$ - The Root Pager}

%%%%%%%%%%%%%%%%%%%%%%%%%%
%%%%%%%%%%%%%%%%%%%%%%%%%%

\label{sec:sigma0}

An \emph{initial address space}, called $\sigma_0$ automatically
exists after the system is booted. $\sigma_0$ maps physical memory
(except for some kernel reserved regions) and is idempotent (all
virtual addresses are the same as the corresponding physical address).
$\sigma_0$ also acts as its own pager\footnote{Note that this is
  somewhat MIPS specific.} by mapping any frame (writable) to the
first task requesting it. Any further requests to $\sigma_0$ for an
already mapped frame is ignored (a null reply is sent).

Figure~\ref{fig:sigma0} illustrates $\sigma_0$ and how it is commonly
used. The inclusion of \(\sigma_0\) in the kernel is specific to
L4/MIPS, \(\sigma_0\) is at user level in other implementations.

One of the first tasks that any \emph{OS personality} needs to perform
is to request all available frames from $\sigma_0$ so that $\sigma_0$
has spent all its memory. This enables the OS to have full control of
memory (other than what is reserved by L4). The OS can then provide
its own pager which maps memory to user tasks (with the appropriate
checks) in response to user page faults.

\begin{figure}[h]
\begin{center}
\leavevmode
\epsfxsize = .7\textwidth
\epsfbox{eps/l4-as.eps}
\end{center}
\caption{$\sigma_0$}
\label{fig:sigma0}
\end{figure}

Part of the kernel reserved space contains the \emph{kernel
  information page} and the \emph{DIT header page}. These special
pages are mapped read-only upon request. Such mappings are part of the
$\sigma_0$ RPC protocol (see Section~\ref{sec:sigma0rpc}).


%%%%%%%%%%%%%%%%%%%%%%%%%%
%%%%%%%%%%%%%%%%%%%%%%%%%%

\subsection{Kernel Information Page}

%%%%%%%%%%%%%%%%%%%%%%%%%%
%%%%%%%%%%%%%%%%%%%%%%%%%%

The kernel information page contains information useful for the
initial servers to find out about the environment they were started
in. Its layout is as follows:

\vspace{1ex}
\begin{center}\begin{tabular}{ll}
\cbox{kernel data}{64}{32} & +40 \\
\cbox{dit header}{64}{32} & +32 \\
\cbox{kernel}{64}{32} & +24 \\
\cbox{memory size}{64}{32} & +16 \\
\cbox{clock}{64}{32} & +8 \\
\cbox{build}{16}{8}\cbox{version}{16}{8}\cbox{``L4uK''}{32}{16} & +0\\
\end{tabular}\end{center}

\vspace{1ex}
\begin{param}{version}
              & L4/R4600 version number.\\
\end{param}
\begin{param}{build}
              & L4/R4600 build number of above version\\
\end{param}
\begin{param}{clock}
              & Number of millisecond ticks since L4 booted.\\
\end{param}
\begin{param}{memory size}
  & The amount, in bytes, of RAM installed on machine L4 is running on.\\
\end{param}
\begin{param}{kernel}
  & The address + 1 of last byte reserved by the kernel of low physical
  memory.\\
\end{param}
\begin{param}{dit header}
  & The address of the DIT header which maps out what was loaded with
  the kernel image.\\
\end{param}
\begin{param}{kernel data}
  & The address of the start of kernel reserved memory in the upper
  physical memory region.
\end{param}


The initial free physical memory available to applications lies
between kernel and kernel data. Figure~\ref{fig:kinfo} shows how the
memory is arranged. Note that the kernel information page and the DIT
header page are within the L4 reserved area while the initial servers
and other data are loaded into the free memory area.

\begin{figure}[h]
\begin{center}
\leavevmode
\epsfxsize = .7\textwidth
\epsfbox{eps/kinfo.eps}
\end{center}
\caption{Physical Memory Arrangement}
\label{fig:kinfo}
\end{figure}

%%%%%%%%%%%%%%%%%%%%%%%%%%
%%%%%%%%%%%%%%%%%%%%%%%%%%

\subsection{DIT}

%%%%%%%%%%%%%%%%%%%%%%%%%%
%%%%%%%%%%%%%%%%%%%%%%%%%%

\label{sec:dit}

DIT is a MIPS specific tool used to build kernel boot images. Like the
L4 kernel itself, it has an information page describing the layout of
various programs and data that were part of the kernel image. It
consists of two parts, the initial header followed by zero of more
file headers as specified by the initial header. The initial headers
format follows below.

\vspace{2ex}
\begin{center}\begin{tabular}{ll}
\cbox{vaddr end}{32}{32} & +20 \\
\cbox{file end}{32}{32} & +16 \\
\cbox{phdr num}{32}{32} & +12 \\
\cbox{phdr size}{32}{32} & +8 \\
\cbox{phdr off}{32}{32} & +4 \\
\cbox{``dhdr''}{32}{32} & +0\\
\end{tabular}\end{center}

\begin{param}{phdr off}
  & The offset from the beginning of this header to where the file
  headers start.\\
\end{param}
\begin{param}{phdr size}
              & The size of each of the  file  headers.\\
\end{param}
\begin{param}{phdr num}
              & The number of file  headers that follow.\\
\end{param}
\begin{param}{file end}
  & The offset to the end of the kernel image file. For DIT internal use
  only.\\
\end{param}
\begin{param}{vaddr end}
  & The end of currently used physical memory space. This includes the
  L4 kernel and all other programs and data in the downloaded kernel
  image.\\
\end{param}

\noindent Each of the program headers is laid out as follows.

\vspace{2ex}
\begin{center}\begin{tabular}{ll}
\cbox{flags}{32}{32} & +28 \\
\cbox{entry}{32}{32} & +24 \\
\cbox{size}{32}{32} & +20 \\
\cbox{base}{32}{32} & +16\\
\cbox{name string}{32}{32} & +12 \\
\cbox{name string}{32}{32} & +8 \\
\cbox{name string}{32}{32} & +4 \\
\cbox{name string}{32}{32} & +0 \\
\end{tabular}\end{center}

\begin{param}{name string}
  & Null terminated string containing name of program or data
  file (truncated to 16 characters).\\
\end{param}
\begin{param}{base}
  & The base address of the program or data file.\\
\end{param}
\begin{param}{size}
  & The size of the program or data.\\
\end{param}
\begin{param}{entry}
              & The start address of the program if its executable, zero
              otherwise.\\
\end{param}
\begin{param}{flags}
              & Miscellaneous flags defined below.\\
\end{param}


\vspace{2ex}
\begin{center}\begin{tabular}{l}
\cbox{\undef}{31}{31}\abox{r}{1}\\
\end{tabular}\end{center}

\begin{param}{r}
  & If set the kernel runs this program as part of the initial servers
  upon startup. If not set, the program or data has simply been loaded
  into memory and has not been invoked.\\
\end{param}

Figure~\ref{fig:dit} shows a more detailed view of the boot image as described by the DIT header. Note that the \emph{files} Section is made up of initial servers followed by other data (possibly code or text).

\begin{figure}[h]
\begin{center}
\leavevmode
\epsfxsize = .7\textwidth
\epsfbox{eps/dit.eps}
\end{center}
\caption{Detailed Boot Image}
\label{fig:dit}
\end{figure}

%%%%%%%%%%%%%%%%%%%%%%%%%%
%%%%%%%%%%%%%%%%%%%%%%%%%%

\subsection{$\sigma_0$ RPC Protocol}

%%%%%%%%%%%%%%%%%%%%%%%%%%
%%%%%%%%%%%%%%%%%%%%%%%%%%

\label{sec:sigma0rpc}

A task can request a mapping from $\sigma_0$ by sending a short
message (see Section~\ref{sec:messages}) to $\sigma_0$. The specific
request is determined by up to the first two words in the register
data of the request. If the request is valid, $\sigma_0$ sends a
mapping to the requester. Note that only tasks in $\sigma_0$'s clan
(i.e.\ the initial servers) can IPC directly to $\sigma_0$ due to the
clans and chiefs mechanisms.

$\sigma_0$ provides four mapping categories:

%%%%%%%%%%%%%%%%%%%%%%%%%%
%%%%%%%%%%%%%%%%%%%%%%%%%%

\subsubsection{Free Physical Memory}

%%%%%%%%%%%%%%%%%%%%%%%%%%
%%%%%%%%%%%%%%%%%%%%%%%%%%

Physical frames can be requested from $\sigma_0$ by sending the
address of the frame as the first register word in a short IPC message
to $\sigma_0$. This is a special form of a page-fault RPC (see
Section~\ref{s:pf}).

\begin{tabular}{ll}
\emph{msg.w0} (\reg{s0}) & \cbox{address}{64}{32}\\
\emph{msg.w1} (\reg{s1}) & \cbox{\undef}{64}{32}\\
\end{tabular}
\vspace{1ex}

If the requested memory is in the available range and not previously
mapped, $\sigma_0$ will send back a writable mapping for that frame to
the requester, otherwise it will send a null reply. \\ 

\example{Requesting physical memory from $\sigma_0$}

This example covers both sending a request to $\sigma_0$ for a
physical frame and successfully receiving the mapping from $\sigma_0$.
The following procedure illustrates explicit memory requests from
$\sigma_0$ but note that the same can be achieved by just touching an
unmapped page (i.e.\ read from or write to the page). The advantage of
just touching a page is decreased complexity but the cost is that
error recovery is not possible: An attempt to obtain a page already
claimed by another initial server would livelock the faulting thread.

\begin{enumerate}

\item Declare a register buffer. \\ \\
  \hspace*{10pt} {\footnotesize\verb+l4_ipc_reg_msg_t msg;+}

\item Declare a variable to hold the address of the desired physical
  frame to be mapped. Assign that address to the variable. \\ \\
  \hspace*{10pt} {\footnotesize\verb+dword_t addr;+} \\
  \hspace*{10pt} {\footnotesize\verb+addr = +} \ldots {\footnotesize\verb+; /* some address */+}
  
\item Copy the frame address to the first word of the register buffer.
  \\ \\  \hspace*{10pt} {\footnotesize\verb+msg.reg[0] = addr;+} \\
  
\item Send the IPC request (in a short message) to $\sigma_0$. \\ \\
  \hspace*{10pt} {\footnotesize\verb+l4_mips_send(SIGMA0_TID, L4_IPC_SHORT_MSG, &msg, L4_IPC_NEVER, &result);+} \\ \\
  {\footnotesize\verb+result+} is the usual result buffer ({\footnotesize\verb+l4_msgdope_t result;+}).
  
\item Assuming all the steps are correct so far, can now wait to
  receive the mapping from $\sigma_0$. \\ \\
  \hspace*{10pt} {\footnotesize\verb+l4_mips_receive(SIGMA0_TID,+}\\
  \hspace*{90pt} {\footnotesize\verb+(void *) ((dword_t) L4_IPC_SHORT_FPAGE |+}\\
  \hspace*{95pt} {\footnotesize\verb+(L4_WHOLE_ADDRESS_SPACE << 2)),+}\\
  \hspace*{90pt} {\footnotesize\verb+&msg, L4_IPC_NEVER, &result);+} \\ \\
  Note that the message reply from $\sigma_0$ is expected to include
  a single page mapping. Recall from the message descriptor discussion
  in Section~\ref{sec:messages} that mappings can be received as part of
  a short message by providing a receive fpage for the message
  descriptor parameter (\emph{rcv\_msg}) instead of the start address
  (or nil) of a message buffer.
  
  The $m$-bit needs to be set to indicate that fpages are to be received.
  This is done via a bit wise \emph{OR} with {\footnotesize\verb+L4_IPC_SHORT_FPAGE+}
  (defined in {\footnotesize\verb+ipc.h+}).
  
  In this example, the mappings are made idempotent by specifying the
  whole address space as the receive fpage. Again, recall that the
  fpage format for this is: \\
  
  \begin{tabular}{ll}
    \hspace*{80pt}&\\[\Up]
                                %
fpage$(0,2^{64} - K)$ &
\cbox{0}{52}{20}\bbox{0}{3}\cbox{64}{7}{7}\abox{\undef{}}{1}\abox{\undef{}}{1}\\[20pt]
%
\end{tabular}

        {\footnotesize\verb+L4_WHOLE_ADDRESS_SPACE+} is defined to be {\footnotesize\verb+64+} (in
        {\footnotesize\verb+types.h+}) and with some thought, it can be seen that
        {\footnotesize\verb+L4_WHOLE_ADDRESS_SPACE << 2+} corresponds exactly to the above
        fpage format. (See the next example for a macro that helps perform
        this task.)
        
\end{enumerate}

For added elegance, it is worth noting that the last two steps above
can be combined by using the appropriate C procedure. An IPC send to a
particular thread followed immediately by a closed receive from the
same thread (usual RPC semantics) is facilitated by the
{\footnotesize\verb+l4_mips_ipc_call+} system call. The last two steps above can then
be replaced by a single step: \\ 

\hspace*{20pt} {\footnotesize\verb+l4_mips_ipc_call(SIGMA0_TID, L4_IPC_SHORT_MSG, &msg,+}\\
\hspace*{110pt} {\footnotesize\verb+(void *) ((dword_t) L4_IPC_SHORT_FPAGE |+}\\
\hspace*{160pt} {\footnotesize\verb+(L4_WHOLE_ADDRESS_SPACE << 2)),+}\\
\hspace*{110pt} {\footnotesize\verb+&msg, L4_IPC_NEVER, &result);+}

%%%%%%%%%%%%%%%%%%%%%%%%%%
%%%%%%%%%%%%%%%%%%%%%%%%%%

\subsubsection{Kernel Information Page}

%%%%%%%%%%%%%%%%%%%%%%%%%%
%%%%%%%%%%%%%%%%%%%%%%%%%%

Unlike free physical frames, the kernel information allows for
multiple mappings to multiple requesters. A request for the kernel
information page can be made by providing a particular invalid address
(-3 on MIPS) as the first register word in a short IPC to $\sigma_0$.
\\ 

\begin{tabular}{ll}
\emph{msg.w0} (\reg{s0}) & \cbox{0xFFFFFFFFFFFFFFFD}{64}{32}\\
\emph{msg.w1} (\reg{s1}) & \cbox{\undef}{64}{32}\\
\end{tabular}
\vspace{1ex}

$\sigma_0$ will respond to a request in the above format by mapping
the kernel information page read-only to the requester. The requester
receives the address of the kernel information page in \textbf{s0}
assuming a one to one mapping.\\ 

\example{Requesting the Kernel Information Page from $\sigma_0$}

This example covers both sending a request to $\sigma_0$ for the
kernel information page and successfully receiving the mapping from
$\sigma_0$.

\begin{enumerate}

\item Declare a register buffer. \\ \\
  \hspace*{10pt} {\footnotesize\verb+l4_ipc_reg_msg_t msg;+}

\item Declare a variable to hold the address of the kernel
  information page (the type {\footnotesize\verb+l4_kernel_info+} is defined in
  {\footnotesize\verb+sigma0.h+}). \\ \\
  \hspace*{10pt} {\footnotesize\verb+l4_kernel_info *kinfo;+}
  
\item Provide the correct invalid address as the first register word.
  The constant {\footnotesize\verb+SIGMA0_KERNEL_INFO_MAP+} defined in {\footnotesize\verb+sigma0.h+}
  can be used for this purpose.
  \\ \\  \hspace*{10pt} {\footnotesize\verb+msg.reg[0] = SIGMA0_KERNEL_INFO_MAP;+} \\
  
\item RPC the request to $\sigma0$. \\ \\
  \hspace*{20pt} {\footnotesize\verb+l4_mips_ipc_call(SIGMA0_TID, L4_IPC_SHORT_MSG, &msg,+}\\
  \hspace*{110pt} {\footnotesize\verb+L4_IPC_MAPMSG(0, L4_WHOLE_ADDRESS_SPACE),+}\\
  \hspace*{110pt} {\footnotesize\verb+&msg, L4_IPC_NEVER, &result);+} \\ \\
  Note that {\footnotesize\verb+L4_IPC_MAPMSG+} is a macro (defined in {\footnotesize\verb+ipc.h+})
  that can be used to construct receive descriptors for short fpage IPCs.
  
\item Store the address of the kernel information page returned by
  $\sigma_0$. \\ \\
  \hspace*{20pt} {\footnotesize\verb+kinfo = (l4_kernel_info *) msg.reg[0];+}

\end{enumerate}

%%%%%%%%%%%%%%%%%%%%%%%%%%
%%%%%%%%%%%%%%%%%%%%%%%%%%

\subsubsection{DIT Header page}

%%%%%%%%%%%%%%%%%%%%%%%%%%
%%%%%%%%%%%%%%%%%%%%%%%%%%

$\sigma_0$ also supports multiple mappings of the DIT header page to
multiple requesters. A request for the DIT header page is made by
providing the DIT header address (obtained from the kernel information
page) as the first register word in a short IPC message to $\sigma_0$.

\begin{tabular}{ll}
\emph{msg.w0} (\reg{s0}) & \cbox{dit header address}{64}{32}\\
\emph{msg.w1} (\reg{s1}) & \cbox{\undef}{64}{32}\\
\end{tabular}
\vspace{1ex}

$\sigma_0$ will respond to a request in the above format by mapping the DIT header page read-only to the requester.  \\

\example{Accessing the DIT Header Page}

This example illustrates how information from the DIT header page can
be obtained. Assuming that the kernel information page has already
been mapped (see previous example) there are two ways to access the
DIT header page:

\begin{itemize}

\item Explicitly request a mapping of the DIT header page from
  $\sigma_0$ before accessing it.

\item Access the DIT header page without an explicit mapping. This
  works because accessing the DIT header page will cause a page fault
  if it has not already been mapped into the task's address space.
  $\sigma_0$ handles the page fault (since the task must be within
  $\sigma_0$'s clan otherwise it can't IPC to $\sigma_0$ anyway) by
  mapping the DIT header page into the faulter's address space.

\end{itemize}

The second method is simpler. The following example uses this method
to access the DIT header page and subsequently each file header.

\begin{enumerate}

\item Declare a variable to hold the address of the DIT header page
  and another variable to hold the address of a file header (the types
  {\footnotesize\verb+Dit_Dhdr+} and {\footnotesize\verb+Dit_Phdr+} are defined in {\footnotesize\verb+dit.h+}).\\ \\
  \hspace*{20pt} {\footnotesize\verb+Dit_Dhr dhdr;+} \\
  \hspace*{20pt} {\footnotesize\verb+Dit_Phdr phdr;+}
  
\item Assign the address of the DIT header page (from the already
  mapped kernel information page - see previous example) to the
  corresponding variable. \\ \\
  \hspace*{20pt} {\footnotesize\verb+dhdr = (Dit_Dhdr *) (kinfo -> dit_hdr);+}
  
\item Assign the address of the first file header to the corresponding
  variable. This address can be calculated by using the file header
  offset found in the DIT header page. \\ \\
  \hspace*{20pt} \verb`phdr = (Dit_Phdr *) ((int) dhdr + dhdr->d_phoff);`
  
\item Cycle through each of the file headers. Recall that the number
  of file headers is also given in the DIT header page. \\ \\
  \hspace*{20pt} \verb`for(i = 0; i < dhdr->d_phnum; i++, phdr++)` \\
  \hspace*{20pt} {\footnotesize\verb+{+} \\
    \hspace*{50pt} {\footnotesize\verb+/* do some work */+} \\
    \hspace*{85pt} \vdots \\
    \hspace*{20pt} {\footnotesize\verb+}+}

\end{enumerate}

%%%%%%%%%%%%%%%%%%%%%%%%%%
%%%%%%%%%%%%%%%%%%%%%%%%%%

\subsubsection{Devices}

%%%%%%%%%%%%%%%%%%%%%%%%%%
%%%%%%%%%%%%%%%%%%%%%%%%%%

A request for a device mapping can be made by providing a particular
invalid address (-2 on MIPS) as the first register word and the device
address as the second register word in a short IPC to $\sigma_0$.
Device addresses are identified by values outside the RAM range.

\begin{tabular}{ll}
\emph{msg.w0} (\reg{s0}) & \cbox{0xFFFFFFFFFFFFFFFE}{64}{32}\\
\emph{msg.w1} (\reg{s1}) & \cbox{address}{64}{32}\\
\end{tabular}
\vspace{1ex}

$\sigma_0$ will respond by mapping \emph{address} writable and
uncacheable to the requester.

%%%%%%%%%%%%%%%%%%%%%%%%%%
%%%%%%%%%%%%%%%%%%%%%%%%%%

\section{\label{s:pf}Page Fault Handling}

%%%%%%%%%%%%%%%%%%%%%%%%%%
%%%%%%%%%%%%%%%%%%%%%%%%%%

\label{sec:pf}

A page fault occurs when a task tries to access (read from or write
to) memory that has not already been mapped into its virtual address
space. A pager is a thread that handles page faults by determining the
actions to be taken in the event of a page fault (usually give a
mapping for the faulting address to the faulter).

When a new task/thread is created, a pager is registered with that
task/thread. The registered pager is then responsible for handling all
the thread's page faults. When a client thread triggers a page fault,
the L4 kernel intercepts the interrupt and sends an RPC to the pager
on the client's behalf. That is, the kernel sends the client's
faulting address and instruction pointer in the first two words of a
short IPC message to the registered pager pretending that the IPC
actually comes from the faulting thread. \\ 

\begin{param}{}
            w0 (\reg{s0})     &  \cbox{fault address / 4}{62}{30}\abox{$w$}{1}\abox{\undef}{1}\\
%
            w1 (\reg{s1})     &  \cbox{faulting user-level IP}{64}{32}\\[12pt]
%
            $w\,\EQ\,0$  &  Read page fault.\\
            $w\,\EQ\,1$  &  Write page fault.
\end{param}

The pager will receive the page fault message as if it were directly
from the faulter. It can then respond by sending an fpage mapping for
the faulting address back to the faulter\footnote{This is standard
  page fault handling but a pager may implement an arbitrary policy,
  in particular, kill any clients that attempt an invalid page
  access.}. The client does not actually receive the pager's mapping
as the mapping is intercepted by the kernel on the client's behalf.
The kernel restarts the client with the new mapping in place.

Figure~\ref{fig:l4-pf} illustrates the RPC paging protocol described. \\

\begin{figure}[h]
\begin{center}
\leavevmode
\epsfxsize = .7\textwidth
\epsfbox{eps/l4-pf.eps}
\end{center}
\caption{Page Fault RPC}
\label{fig:l4-pf}
\end{figure}

$\sigma_0$ (see Section~\ref{sec:sigma0}) is an example of a pager
service. It is in fact the root pager because it handles the initial
address space and all mappings can be traced back to $\sigma_0$. The
other type of pagers which stem from $\sigma_0$ are the \emph{external
  pagers}. External pagers are user level threads which perform the
task of page fault handling for other threads. External pagers
themselves originally obtain their mappings from either $\sigma_0$ or
from an intermediate pager (which is just another external pager - see
Figure~\ref{fig:sigma0}).

A pager may send more than one page at a time, this allows implementing
pre-paging.

The \micro-kernel will ignore any attempts to map a page over an
already existing mapping, any existing mapping must first be removed
(e.g.\ by using an {\footnotesize\verb+le_fpage_unmap+} system call). An exception is
an attempt to replace an existing mapping by a mapping to the same
frame. In this case, the write-protection attribute of the mapping is
taken from the present IPC operation (the \(w\)-bit in the send fpage,
see Section~\ref{s:w-bit}). This allows the pager to change the
mapping between read-only and writable.

\example{A Simple Pager}

This example presents the basic tasks that a simple user level pager
needs to perform.

\begin{enumerate}

\item Assume the following declarations. \\ \\
  \hspace*{20pt} {\footnotesize\verb+int r;+}\\
  \hspace*{20pt} {\footnotesize\verb+l4_ipc_reg_msg_t msg;+}\\
  \hspace*{20pt} {\footnotesize\verb+l4_threadid_t thrdid;+}\\
  \hspace*{20pt} {\footnotesize\verb+l4_msgdope_t result;+}\\
  \hspace*{20pt} {\footnotesize\verb+dword_t fault_addr;+}\\
  \hspace*{20pt} {\footnotesize\verb+l4_snd_fpage_t *fp;+}

\item First, the pager needs to wait for a page fault message from
  any client. \\ \\
  \hspace*{20pt} {\footnotesize\verb+r = l4_mips_ipc_wait(&thrdid, L4_IPC_SHORT_MSG, &msg,+}\\
  \hspace*{130pt} {\footnotesize\verb+L4_IPC_NEVER, &result);+} \\ \\
  The faulter is recorded as the sender of the message (in {\footnotesize\verb+thrdid+})
  and the faulting address in the first word of the register buffer.
  
\item The pager now needs to obtain a mapping of the faulting
  address for itself (if it does not already have it) before it can
  pass the mapping on to the faulter. For this example, assume the pager
  gets its mappings from $\sigma_0$. \\ \\
  \hspace*{20pt} {\footnotesize\verb+fault_addr = (msg.reg[0] & (~(dword_t) 3)); /* mask out lower bits */+}\\
  \hspace*{20pt} {\footnotesize\verb+msg.reg[0] = fault_addr | 2; /* request write mapping */+}\\ \\
  \hspace*{20pt} {\footnotesize\verb+r = l4_mips_ipc_call(SIGMA0_TID, L4_IPC_SHORT_MSG, &msg,+}\\
  \hspace*{130pt} {\footnotesize\verb+L4_IPC_MAPMSG(0, L4_WHOLE_ADDRESS_SPACE),+}\\
  \hspace*{130pt} {\footnotesize\verb+&msg, L4_IPC_NEVER, &result);+}
  
\item Construct the send fpage descriptor for the  mapping to the
  faulting client. The single mapping to be sent is the same size as the
  single hardware page ($s$ = 12) and has the fault address as the send
  fpage base ($b$ = fault address) and also the hot spot. Note that the
  function {\footnotesize\verb+l4_fpage+} (defined in {\footnotesize\verb+types.h+}) can be used to
  construct an fpage. \\ \\
  \hspace*{20pt} {\footnotesize\verb+fp = (l4_snd_fpage_t *) &msg.reg[0];+} \\
  \hspace*{20pt} {\footnotesize\verb+fp[0].snd_base = fault_addr;+}\\
  \hspace*{20pt} {\footnotesize\verb+fp[0].fpage = l4_fpage(fault_addr, 12, 1, 0);+}\\
  \hspace*{20pt} {\footnotesize\verb+fp[1].fpage.fpage = 0; /* invalid fpage as a terminator */+}
  
\item Finally, send the mapping to the client. \\ \\
  \hspace*{20pt} {\footnotesize\verb+r = l4_mips_ipc_send(thrdid, L4_IPC_SHORT_FPAGE, &msg,+}\\
  \hspace*{130pt} {\footnotesize\verb+L4_IPC_NEVER, &result);+}
  
\item The pager returns to the start to wait for the next client
  page fault.

\end{enumerate}

%%%%%%%%%%%%%%%%%%%%%%%%%%
%%%%%%%%%%%%%%%%%%%%%%%%%%

\section{Exception And Interrupt Handling}

%%%%%%%%%%%%%%%%%%%%%%%%%%
%%%%%%%%%%%%%%%%%%%%%%%%%%

%%%%%%%%%%%%%%%%%%%%%%%%%%
%%%%%%%%%%%%%%%%%%%%%%%%%%

\subsection{Exception Handling}

%%%%%%%%%%%%%%%%%%%%%%%%%%
%%%%%%%%%%%%%%%%%%%%%%%%%%

Exceptions are handled in a similar fashion to page faults (on
L4/MIPS). As with registering pagers, when a new task/thread is
created, an exception handler thread is also registered with it. When
a thread raises an exception, the kernel's interrupt handler sends an
RPC to the thread's registered exception handler. The protocol is
similar to that depicted in Figure~\ref{fig:l4-pf}. The exception
handler can choose to continue the thread that raised the exception by
replying to the message. Ignoring the IPC will block the faulting
thread forever and thus effectively kill it.

%%%%%%%%%%%%%%%%%%%%%%%%%%
%%%%%%%%%%%%%%%%%%%%%%%%%%

\subsection{Interrupt Handling}

%%%%%%%%%%%%%%%%%%%%%%%%%%
%%%%%%%%%%%%%%%%%%%%%%%%%%

\label{sec:interrupts}

A thread can install itself as an interrupt handler for a particular
interrupt by associating itself with the interrupt. A thread that
wants to associate itself with an interrupt needs to invoke a
\emph{receive} IPC specifying the interrupt number as the sender
(\emph{src} parameter) and a zero timeout. That is: \\ 

\hspace*{20pt} {\footnotesize\verb+l4_threadid_t sender;+} \\
\hspace*{20pt} {\footnotesize\verb+l4_ipc_reg_msg_t msg;+} \\
\hspace*{20pt} {\footnotesize\verb+l4_msgdope_t result;+}\\ \\
\hspace*{20pt} {\footnotesize\verb+sender = +} \ldots {\footnotesize\verb+; /* interrupt number */+}\\
\hspace*{20pt} {\footnotesize\verb+l4_mips_ipc_receive(sender, L4_IPC_SHORT_MSG, &msg,+}\\
\hspace*{126pt} {\footnotesize\verb+L4_IPC_TIMEOUT(0, 0, 0, 1, 0, 0), &result);+}\\

Each interrupt is assigned to the first thread that attempts to
associate with it. Any attempt to associate with an already associated
interrupt will fail. A thread can dissociate itself as a handler for
an interrupt by associating with a NULL interrupt.

Once associated with an interrupt, a handler waits for interrupts by
receiving with non-zero timeout.

%%%%%%%%%%%%%%%%%%%%%%%%%%
%%%%%%%%%%%%%%%%%%%%%%%%%%

\section{OS On L4}

%%%%%%%%%%%%%%%%%%%%%%%%%%
%%%%%%%%%%%%%%%%%%%%%%%%%%

This section discusses some of the basic issues and concepts that need
to be considered when creating an OS with L4 as the base.

%%%%%%%%%%%%%%%%%%%%%%%%%%
%%%%%%%%%%%%%%%%%%%%%%%%%%

\subsection{\label{s:u:osstruct}OS Structure}

%%%%%%%%%%%%%%%%%%%%%%%%%%
%%%%%%%%%%%%%%%%%%%%%%%%%%

The \micro-kernel concept logically leans towards a client-server OS
model. The OS provides servers that perform the required operations
upon request from clients (the pager is an example of this - see
Section~\ref{sec:pf}). Each OS server and the entire OS itself are in
fact operating at user level (same as the client). As such, multiple
\emph{OS personalities} can exist simultaneously without affecting
each other. Figure~\ref{fig:l4-os} shows a typical L4-based OS
structure.\\ 

\begin{figure}[h]
\begin{center}
\leavevmode
\epsfxsize = .7\textwidth
\epsfbox{eps/l4-os.eps}
\end{center}
\caption{Typical L4-based OS Structure}
\label{fig:l4-os}
\end{figure}

In L4 terms, an OS may consist of many server threads in the same or
separate tasks (i.e. single or multiple task OS). One of the OS tasks
will start up client processes making the OS chief of each of the
client processes.  Clients then make ``system calls'' via library
stubs which actually perform RPC to an OS server. The OS server may
choose to service the request itself or can redirect the request to
another OS server within or outside of its own task.

For example, the OS may contain a system call handler server which
makes use of RPC redirection.  The steps involved in a single client
system call would then be:

\begin{enumerate}

\item The client uses the library stub to send an RPC to the system
  call server and waits for a reply.

\item System call server forwards the request to another server via a
  \emph{deceiving} send IPC operation giving the client id as the
  virtual sender.

\item The subsystem performs the necessary operation and replies
  directly back to the client via a \emph{deceiving} send IPC
  operation giving the system call id as the virtual sender.

\item The client receives the reply to the original request as if it
  came directly from the system call server.

\end{enumerate}

Steps 2 and 3 above uses deceiving IPC to maintain RPC semantics.
Figure~\ref{fig:l4-redir} illustrates this RPC redirection procedure.
\\ 

\begin{figure}[h]
\begin{center}
\leavevmode
\epsfxsize = .5\textwidth
\epsfbox{eps/l4-redir.eps}
\end{center}
\caption{RPC Redirection}
\label{fig:l4-redir}
\end{figure}

%%%%%%%%%%%%%%%%%%%%%%%%%%
%%%%%%%%%%%%%%%%%%%%%%%%%%

\subsection{Some Conventions}

%%%%%%%%%%%%%%%%%%%%%%%%%%
%%%%%%%%%%%%%%%%%%%%%%%%%%

%%%%%%%%%%%%%%%%%%%%%%%%%%
%%%%%%%%%%%%%%%%%%%%%%%%%%

\subsubsection{OS Startup}

%%%%%%%%%%%%%%%%%%%%%%%%%%
%%%%%%%%%%%%%%%%%%%%%%%%%%

When the OS first starts, there are a few recommended tasks that it
should perform before starting any client processes. Basically, the OS
needs to acquire all resources which are allocated on a first request
basis (to ensure resource control and prevent loss of resources to
clients) as well as set up the necessary structures. The tasks the OS
should perform on startup include:

\begin{itemize}

\item Register for all free interrupts.

\item Request all free memory from $\sigma_0$.

\item Request all pages containing device addresses (unless these are
  requested by separate driver tasks, which must then also be initial
  servers)

\item Set up data structures for memory management (reserve space for
  free lists, frame table, client page tables etc.)

\item Acquire all inactive tasks (i.e.\ use {\footnotesize\verb+l4_task_new+} on all
  inactive task IDs).

\item Start other server threads (if any).

\item Start device driver threads (if any).

\item Set up service data structures (i.e. TCBs, file system, etc.).

\end{itemize}

After the OS has completed all the preliminary tasks, it can start up
one or many client tasks. The OS may choose to donate some of its
acquired tasks to any of the user (sub)tasks.

Note that the OS itself is just a user level L4 task and thus can be
interrupted and scheduled out. This fact, together with
multi-threading, makes it essential to have concurrency control on all
OS data structures.

Initial servers are started up by L4 and therefore cannot receive a
stack address from their parent. They must therefore initialise their
own stack. This is normally done by using special startup code for
initial servers (``crtS'') which initialises the stack pointer to the
top of a static array.

%%%%%%%%%%%%%%%%%%%%%%%%%%
%%%%%%%%%%%%%%%%%%%%%%%%%%

\subsubsection{Clients}

%%%%%%%%%%%%%%%%%%%%%%%%%%
%%%%%%%%%%%%%%%%%%%%%%%%%%

After start up, the OS will be either waiting for a client request or
actually handling a client request. After the preliminaries, it is the
job of the OS to start the client tasks (else the OS will have nothing
to do!). But where does the OS look to find clients to start? The
answer is that it is a matter of convention and is decided by the OS
designer. Some common options are:

\begin{itemize}

\item OS starts up first non-initial server executable in the boot
  image. (Recall that the files section of the boot image is composed
  of initial servers followed by other data - see
  Figure~\ref{fig:dit}.  Non-initial servers are part of the other
  data).

\item First non-executable item in the boot image contains a list of
  initial client tasks.

\item The OS looks for an executable with a predefined name in the
  boot image.

\end{itemize}

After the client has been started, it will want to make system calls
as a normal part of its operations. Again, there needs to be some
convention defined to let clients know where to send their system call
RPCs (i.e.\ need to know id of OS's system call server). Some common
options are:

\begin{itemize}

\item First thing the client does is an open receive. The OS (or any
  other task that started the client) follows the protocol by sending
  a message to the client containing its own identity or the identity
  of another server. Note that this client startup protocol can be
  hidden in startup code (``crt0'').

\item Client could send all system call RPCs to $\sigma_0$. Then the
  clans and chiefs mechanism will ensure that the chief of the client
  task will intercept the message. The chief can then reply pretending
  to be $\sigma_0$ via a deceiving send IPC.\\ 

\end{itemize}

\example{Client Startup Protocol}

This example illustrates one way in which the OS can disclose the
identity of its system call server id to client tasks. \\ \\ 

\hspace*{10pt} OS code fragment: \\ \\
\hspace*{50pt} \vdots \\
\hspace*{20pt} {\footnotesize\verb+l4_ipc_reg_msg_t msg;+} \\
\hspace*{20pt} {\footnotesize\verb+l4_msg_dope_t result;+} \\
\hspace*{20pt} {\footnotesize\verb+l4_threadid_t subid, syscallid;+} \\ \\
\hspace*{20pt} {\footnotesize\verb+syscallid = +} \ldots {\footnotesize\verb+; /* system call server id */+} \\
\hspace*{20pt} {\footnotesize\verb+subid = l4_task_new(+} \ldots {\footnotesize\verb+); /* start the client task */+} \\ \\
\hspace*{30pt} {\footnotesize\verb+ /* Send system call server id to the new client */+} \\
\hspace*{20pt} {\footnotesize\verb+ msg.reg[0] = syscallid.ID;+} \\
\hspace*{20pt} {\footnotesize\verb+ l4_mips_ipc_send(subid, L4_IPC_SHORT_MSG, &msg, L4_IPC_NEVER, &result);+} \\
\hspace*{50pt} \vdots \\ \\

\hspace*{10pt} Client code fragment: \\ \\
\hspace*{20pt} {\footnotesize\verb+l4_ipc_reg_msg_t msg;+} \\
\hspace*{20pt} {\footnotesize\verb+l4_msgdope_t result;+} \\
\hspace*{20pt} {\footnotesize\verb+l4_threadid_t OSid, syscallid;+} \\ \\
\hspace*{30pt} {\footnotesize\verb+/* Wait for message from OS */+} \\
\hspace*{20pt} {\footnotesize\verb+l4_mips_ipc_wait(&OSid, L4_IPC_SHORT_MSG, &msg, L4_IPC_NEVER, &result);+} \\ \\
\hspace*{30pt} {\footnotesize\verb+/* Store system call server id from message */+} \\
\hspace*{20pt} {\footnotesize\verb+syscallid = (l4_threadid_t) msg.reg[0];+} \\
\hspace*{50pt} \vdots \\

Note that the client gets both the id of the OS (implicitly - as the
sender in the open receive) and the system call server id (explicitly
- as part of the message).

% 
%%% Local Variables: %%%
%%% mode: latex %%%
%%% End: %%%


%\include{bib}
\bibliographystyle{alpha}
\bibliography{defs,os,dist}

%%%%%%%%%%%%%%%%%%%%%%%%%%%%%%%%%%%%%%%%%%%%%%%%%%%%%%%%%%%%%%%%%%%%%%

\end{document}

% 
%%% Local Variables: %%%
%%% mode: latex %%%
%%% End: %%%
